\documentclass[12pt,a4paper,oneside]{report}

\usepackage[a4paper,margin=2.5cm]{geometry} % SEU: normal margins, approximately 1 inch
\usepackage{fontspec}
\usepackage{polyglossia}

% ========================================
% ენის და შრიფტის კონფიგურაცია (SEU)
% ========================================
\setdefaultlanguage{georgian} % ქართული, როგორც მთავარი ენა
\setotherlanguage{english}

% Sylfaen არის SEU-ს მიერ მოთხოვნილი ოფიციალური შრიფტი
\setmainfont[
  Ligatures=TeX,
  SizeFeatures={Size=12}
]{Sylfaen}

\newfontfamily\georgianfont[
  SizeFeatures={Size=12}
]{Sylfaen}
\newfontfamily\georgianfonttt[
  SizeFeatures={Size=12}
]{Sylfaen}

\usepackage{setspace}
\usepackage{microtype}
\usepackage{graphicx}
\usepackage{booktabs}
\usepackage{longtable}
\usepackage{array}
\usepackage{multirow}
\usepackage{amsmath,amssymb,amsthm}
\usepackage{siunitx}
\usepackage{enumitem}
\usepackage{float}
\usepackage{caption}
\usepackage{subcaption}
\usepackage{hyperref}
\usepackage{fancyhdr}
\usepackage{titlesec}

\hypersetup{
  colorlinks=true,
  linkcolor=blue,
  citecolor=blue,
  urlcolor=blue,
  pdftitle={CXR-Synapse: გულ-მკერდის რენტგენოგრამების მრავალლეიბლიანი კლასიფიკაცია LISA, PaQ და კონფორმალური უსაფრთხოების გარანტიებით},
  pdfauthor={ნიკა}
}

% ========================================
% ფორმატირების მოთხოვნები (SEU)
% ========================================
\setstretch{1.5} % 1.5 ხაზთაშორისი დაშორება
\setlength{\parindent}{1.2em}
\setlength{\parskip}{0.2em}

% სქოლიოს ფორმატირება: ზომა 10, ერთმაგი დაშორება
\usepackage[perpage,bottom]{footmisc}
\renewcommand{\footnotelayout}{\fontsize{10}{10}\selectfont}
\renewcommand{\footnoterule}{\kern -3pt \hrule \kern 2.6pt}
\renewcommand{\thefootnote}{\arabic{footnote}}

\pagestyle{fancy}
\setlength{\headheight}{15pt}
\fancyhf{}
\fancyhead[R]{\thepage}
\fancyhead[L]{სამაგისტრო ნაშრომი - CXR-Synapse}
\renewcommand{\headrulewidth}{0.4pt}

% თავების სათაურები: 16pt, მუქი (Bold), ცენტრირებული
\titleformat{\chapter}[display]
  {\bfseries\fontsize{16}{19}\selectfont\filcenter}
  {\chaptertitlename\ \thechapter}
  {0.5em}
  {\fontsize{16}{19}\selectfont}

\titleformat{name=\chapter,numberless}[display]
  {\bfseries\fontsize{16}{19}\selectfont\filcenter}
  {}
  {0pt}
  {\fontsize{16}{19}\selectfont}

% ქვეთავების სათაურები: 14pt, მუქი (Bold)
\titleformat{\section}[hang]
  {\bfseries\fontsize{14}{16}\selectfont}
  {\thesection}
  {1em}
  {}

\titleformat{\subsection}[hang]
  {\bfseries\fontsize{12}{14}\selectfont}
  {\thesubsection}
  {1em}
  {}

% ცვლადების დეკლარაცია თავფურცლისთვის
\newcommand{\UniversityName}{საქართველოს ეროვნული უნივერსიტეტი სეუ}
\newcommand{\ProgramName}{სამაგისტრო პროგრამა: ხელოვნური ინტელექტი}
\newcommand{\ThesisTitle}{CXR-Synapse: გულ-მკერდის რენტგენოგრამების მრავალლეიბლიანი კლასიფიკაცია LISA, PaQ და კონფორმალური უსაფრთხოების გარანტიებით}
\newcommand{\AuthorName}{ნიკა [მიუთითეთ გვარი]}
\newcommand{\SupervisorName}{[ხელმძღვანელის სახელი და გვარი]}
\newcommand{\SupervisorDegree}{PhD}
\newcommand{\DegreeStatement}{სამაგისტრო ნაშრომი შესრულებულია ხელოვნური ინტელექტის მაგისტრის აკადემიური ხარისხის მოსაპოვებლად}
\newcommand{\SubmissionYear}{2026}
\newcommand{\SubmissionCity}{თბილისი}

\begin{document}

\pagenumbering{gobble}

% -------------------------
% თავფურცელი (SEU-ს სტანდარტით)
% -------------------------
\begin{titlepage}
    \centering
    {\fontsize{18}{22}\selectfont \UniversityName\par}
    \vspace{3\baselineskip}

    {\fontsize{16}{20}\selectfont \AuthorName\par}
    \vspace{3\baselineskip}

    {\fontsize{16}{20}\selectfont \bfseries \ThesisTitle\par}
    \vspace{3\baselineskip}

    {\fontsize{14}{18}\selectfont \ProgramName\par}
    \vspace{0.7\baselineskip}
    {\fontsize{14}{18}\selectfont \DegreeStatement\par}

    \vfill
    \begin{flushright}
        {\fontsize{14}{18}\selectfont ხელმძღვანელი: \SupervisorName\\
        აკადემიური ხარისხი: \SupervisorDegree}
    \end{flushright}

    \vfill
    {\fontsize{14}{18}\selectfont \SubmissionCity, \SubmissionYear\par}
\end{titlepage}

% -------------------------
% აკადემიური პატიოსნების დეკლარაცია
% -------------------------
\chapter*{აკადემიური პატიოსნების დეკლარაცია}
\addcontentsline{toc}{chapter}{აკადემიური პატიოსნების დეკლარაცია}
ვაცხადებ, რომ წინამდებარე სამაგისტრო ნაშრომი წარმოადგენს ჩემს დამოუკიდებელ ნაშრომს და მასში გამოყენებული ყველა წყარო სათანადოდ არის მითითებული და ციტირებული აკადემიური სტანდარტების დაცვით. 
ჩემს მიერ წარმოდგენილი ექსპერიმენტული მონაცემები, სტატისტიკური ანალიზი და პროგრამული კოდი ავთენტურია. ნაშრომში არ არის გამოყენებული პლაგიატი ან მონაცემთა ფალსიფიკაცია, რაც არღვევს უნივერსიტეტის აკადემიური ეთიკის პოლიტიკას.

\vspace{1.2cm}
\noindent ხელმოწერა: \rule{6cm}{0.4pt} \hfill თარიღი: \rule{3.5cm}{0.4pt}

% -------------------------
% ქართული ანოტაცია
% -------------------------
\chapter*{ანოტაცია}
\addcontentsline{toc}{chapter}{ანოტაცია}
გულ-მკერდის რენტგენოგრამების (CXR) მრავალლეიბლიანი კლასიფიკაცია თანამედროვე სამედიცინო ხელოვნური ინტელექტის ერთ-ერთ მთავარ გამოწვევას წარმოადგენს. ძირითად პრობლემებს შორისაა კლასთა უკიდურესი დისბალანსი, წარმოდგენების დრიფტი (Representational Drift) და გაყინული ვიზუალური მოდელების ზე-თავდაჯერებულობა (Overconfidence). წინამდებარე ნაშრომში წარმოდგენილია \textbf{CXR-Synapse} — კლინიკურად ორიენტირებული არქიტექტურა, რომელიც პირველად აერთიანებს LISA (Language-Invariant Semantic Anchoring) რეგულარიზაციას, PaQ (Pathology-as-Query) სივრცით ყურადღების მექანიზმს 2D ანატომიური პოზიციონირებით და ბეტა-ევიდენციურ ლოგიკას.

ჩატარებულმა ექსპერიმენტებმა ChestMNIST მონაცემთა ბაზაზე აჩვენა, რომ შემოთავაზებული მეთოდი არღვევს ტრადიციულ ზღვარს და აღწევს მრავალლეიბლიან AUC 0.7913-ს. უმთავრესი მიღწევაა მოდელის სრულყოფილი კალიბრაცია, რისი დადასტურებაცაა რეკორდულად დაბალი შეცდომის მაჩვენებელი — ECE 0.0120. ნაშრომში ასევე ინტეგრირებულია Micro-Marginal კონფორმალური რისკის კონტროლი, რაც უზრუნველყოფს იშვიათი პათოლოგიების 92.4\%-იან კლინიკურ დაფარვას და აბათილებს მოდელის მიერ ცრუ პროგნოზების გენერირების რისკს. 

აღნიშნული შედეგები და სტატისტიკურად სარწმუნო (p=0.000) არქიტექტურა ქმნის მყარ საფუძველს Trustworthy AI სისტემების კლინიკურ პრაქტიკაში უსაფრთხოდ დასანერგად.

\noindent\textbf{საკვანძო სიტყვები:}
გულ-მკერდის რენტგენოგრაფია, მრავალლეიბლიანი კლასიფიკაცია, Pathology-as-Query, LISA, კალიბრაცია, ევიდენციური ღრმა სწავლება, კონფორმალური უსაფრთხოება.

% -------------------------
% ინგლისური აბსტრაქტი
% -------------------------
\begin{english}
\chapter*{Abstract}
\addcontentsline{toc}{chapter}{Abstract (in English)}

Chest X-ray (CXR) multi-label classification remains challenging due to severe class imbalance, representational drift, and model overconfidence in frozen vision backbones. This thesis proposes \textbf{CXR-Synapse}, a clinically grounded, state-of-the-art framework that integrates LISA (Language-Invariant Semantic Anchoring), PaQ (Pathology-as-Query) spatial cross-attention with 2D anatomical encodings, and Strictly Proper Beta Evidential logic for robust uncertainty quantification.

Experimental results on the ChestMNIST dataset demonstrate that the proposed architecture overcomes traditional representational bottlenecks, achieving a multi-label AUC of 0.7913. The paramount scientific contribution is the near-perfect probabilistic calibration of the model, yielding an Expected Calibration Error (ECE) of 0.0120. Furthermore, the thesis implements Micro-Marginal Conformal Risk Control, successfully guaranteeing a 92.4\% clinical coverage rate while explicitly balancing the fundamental mathematical tradeoff between discriminative power and prediction set size.

The statistically significant findings (p=0.000) and the rigorous, publication-level evaluation protocol position this framework as a defensible and highly reliable blueprint for deploying Trustworthy AI in high-stakes clinical environments.

\noindent\textbf{Keywords:}
Chest X-ray, multi-label classification, Pathology-as-Query, LISA, calibration, Evidential Deep Learning, Conformal Prediction.
\end{english}

% -------------------------
% სარჩევი, ფიგურები, ცხრილები
% -------------------------
\tableofcontents
\clearpage
\pagenumbering{arabic}
\setcounter{page}{1}

\listoffigures
\listoftables
\clearpage

\chapter*{აბრევიატურების ჩამონათვალი}
\addcontentsline{toc}{chapter}{აბრევიატურების ჩამონათვალი}
\begin{longtable}{p{0.18\textwidth}p{0.76\textwidth}}
\textbf{AI} & Artificial Intelligence (ხელოვნური ინტელექტი) \\
\textbf{AUC} & Area Under the Receiver Operating Characteristic Curve (ROC მრუდქვეშა ფართობი) \\
\textbf{BCE} & Binary Cross-Entropy (ბინარული ჯვარედინი ენტროპია) \\
\textbf{CI} & Confidence Interval (ნდობითი ინტერვალი) \\
\textbf{CXR} & Chest X-Ray (გულ-მკერდის რენტგენოგრამა) \\
\textbf{ECE} & Expected Calibration Error (კალიბრაციის მოსალოდნელი შეცდომა) \\
\textbf{LISA} & Language-Invariant Semantic Anchoring \\
\textbf{PaQ} & Pathology-as-Query \\
\textbf{SOTA} & State of the Art (მიმდინარე ტექნოლოგიური პიკი) \\
\textbf{SWA} & Stochastic Weight Averaging \\
\textbf{VLM} & Vision-Language Model (ვიზუალურ-ტექსტური მოდელი) \\
\end{longtable}

\chapter{შესავალი}

\section{პრობლემის აქტუალობა}
გულ-მკერდის რენტგენოგრაფია პრაქტიკაში ყველაზე ხშირად გამოყენებული დიაგნოსტიკური მეთოდია,
რის გამოც მისი ავტომატიზებული ანალიზი განსაკუთრებულ მნიშვნელობას იძენს კლინიკური ტრაიაჟის,
მუშაობის სიჩქარისა და ხარისხის გაუმჯობესებისთვის.
მიუხედავად ამისა, მრავალლეიბლიანი კლასიფიკაცია კვლავ რთულ ამოცანად რჩება,
რადგან ერთდროულად გვხვდება კლასთა დისბალანსი, პათოლოგიათა თანხვედრა,
ლეიბლების არასრულყოფილი ან არასტაბილური ნიშნვა და მონაცემთა განაწილების ცვლა სხვადასხვა კლინიკურ წყაროს შორის.
შესაბამისად, მხოლოდ მაღალი AUC მაჩვენებელი საკმარისი არ არის;
აუცილებელია სანდოობის, კალიბრაციისა და გაურკვევლობის კონტროლირებადი შეფასებაც.

\section{კვლევის მიზანი}
კვლევის მიზანია შეიქმნას და შეფასდეს CXR-Synapse ტიპის ერთიანი მეთოდოლოგიური ჩარჩო,
რომელიც ELIXR-C/CXR-Foundation ვიზუალურ რეპრეზენტაციებს აერთიანებს RadLex-ით განპირობებულ
ლეიბლ-გრაფულ მიზეზობრიობასთან, long-tail ოპტიმიზაციასთან და uncertainty-aware შეფასებასთან,
რათა გაიზარდოს როგორც დიაგნოსტიკური ხარისხი, ისე პროგნოზის სანდოობა კლინიკურად რელევანტურ პირობებში.

\section{კვლევის ამოცანები}
დასმული მიზნის მისაღწევად განისაზღვრა შემდეგი ამოცანები:
\begin{enumerate}[leftmargin=1.6em]
    \item თანამედროვე მრავალლეიბლიანი CXR მეთოდებისა და ბენჩმარკების ანალიზი;
    \item ELIXR-C/CXR-Foundation-ზე დაფუძნებული ვიზუალური ნიშნების ინტეგრაცია ონტოლოგიურად განპირობებულ ლეიბლ-გრაფთან;
    \item კლასთა დისბალანსზე ორიენტირებული (long-tail-aware) დანაკარგების ფუნქციის შერჩევა და დასაბუთება;
    \item გაურკვევლობის შეფასების მექანიზმის აგება deep ensemble მიდგომით;
    \item კალიბრაციისა (ECE, Brier) და კონფორმალური პროგნოზის (coverage/set size) შეფასება;
    \item შედეგების სტატისტიკური სანდოობის დადგენა bootstrap ნდობითი ინტერვალებით და დაწყვილებული ტესტებით.
\end{enumerate}

\section{კვლევის მეთოდოლოგია}
კვლევის მეთოდოლოგია შერჩეულია ისე, რომ თითოეული მეთოდი პირდაპირ პასუხობდეს მიზანსა და ამოცანებს,
და უზრუნველყოფდეს მიღებული დასკვნების სანდოობას.

\textbf{1) შედარებითი სამეცნიერო ანალიზი.}
აქტუალური ლიტერატურის სისტემური მიმოხილვა გამოიყენება პრობლემის ფორმალიზაციისა და
არსებული მეთოდების ძლიერი/სუსტი მხარეების განსაზღვრისთვის.
ეს მეთოდი რელევანტურია, რადგან იძლევა კვლევის ხარვეზების იდენტიფიცირების საშუალებას და
ახალი მეთოდის პოზიციონირებას არსებულ მეცნიერებრივ კონტექსტში.

\textbf{2) ექსპერიმენტული ღრმა სწავლება კონტროლირებადი აბლაციებით.}
მოდელის კომპონენტების წვლილი იზომება აბლაციური ექსპერიმენტებით,
რაც უზრუნველყოფს მიზეზ-შედეგობრივი ინტერპრეტაციის შესაძლებლობას.
რელევანტურობა განპირობებულია იმით, რომ სწორედ აბლაციები აჩვენებს
რომელი მოდული ქმნის რეალურ გაუმჯობესებას.

\textbf{3) სტატისტიკური დასკვნითი ანალიზი.}
bootstrap ნდობითი ინტერვალები და დაწყვილებული მნიშვნელოვნების ტესტები გამოიყენება
შედეგების შემთხვევითი ვარიაციისგან გასამიჯნად.
მეთოდის სანდოობა ეფუძნება განმეორებად ნიმუშებზე შეფასებას და
შედეგების რაოდენობრივად დასაბუთებას.

\textbf{4) კალიბრაციისა და გაურკვევლობის შეფასება.}
ECE/Brier, deep ensemble variance და conformal coverage-set size მეტრიკები
გამოიყენება პროგნოზის სანდოობის შესაფასებლად.
ეს მეთოდები პასუხობს პრაქტიკულ მიზანს: სისტემამ არა მხოლოდ სწორად უნდა კლასიფიციროს,
არამედ უნდა წარმოადგინოს სანდოობის ინტერპრეტირებადი ხარისხი.

\textbf{5) ხარისხობრივი შეცდომების ტაქსონომია.}
ცრუ-დადებითი/ცრუ-უარყოფითი შემთხვევების სტრუქტურირებული ანალიზი
უზრუნველყოფს მეთოდის გამოყენებადობის კრიტიკულ შეფასებას კლინიკურ გარემოში.

ზემოაღნიშნული მეთოდების ერთობლიობა უზრუნველყოფს
(ა) კვლევის მიზნებთან შესაბამისობას,
(ბ) შედეგების შიდა და გარე ვალიდობის ზრდას,
და (გ) ნაშრომის რეპროდუცირებადობის მაღალ დონეს.

\section{ლიტერატურის მიმოხილვა}
კვლევის თეორიულ საფუძველს ქმნის რამდენიმე მიმართულება:
(1) დიდი საჯარო CXR ბენჩმარკები (ChestX-ray14, CheXpert, MIMIC-CXR, MedMNIST v2)
\cite{wang2017chestxray14,irvin2019chexpert,johnson2019mimiccxr,yang2023medmnistv2};
(2) თვითზედამხედველური ვიზუალური რეპრეზენტაციები DINO/MAE ოჯახიდან
\cite{caron2021dino,he2022mae,oquab2024dinofamily};
(3) მრავალლეიბლიანი დისბალანსური კლასიფიკაციის დანაკარგების ფუნქციები
\cite{ridnik2021asymmetric,leng2022polyloss};
(4) გაურკვევლობისა და კალიბრაციის თანამედროვე ჩარჩოები
\cite{lakshminarayanan2017ensembles,guo2017calibration,angelopoulos2021conformal}.
ლიტერატურის დეტალური, სისტემატიზებული განხილვა წარმოდგენილია ნაშრომის მეორე თავში.

\section{ნაშრომის აგებულების მოკლე აღწერა}
ნაშრომის პირველი თავი ეთმობა პრობლემის აქტუალობის, მიზნის, ამოცანებისა და მეთოდოლოგიის დასაბუთებას.
მეორე თავში წარმოდგენილია თემასთან დაკავშირებული თანამედროვე სამეცნიერო ლიტერატურის მიმოხილვა.
მესამე თავში აღწერილია CXR-Synapse მეთოდის არქიტექტურა და ოპტიმიზაციის ჩარჩო.
მეოთხე თავში განსაზღვრულია ექსპერიმენტული პროტოკოლი, რეპროდუცირებადობის წესები და სტატისტიკური შეფასების მიდგომები.
მეხუთე თავში წარმოდგენილია რაოდენობრივი და ხარისხობრივი შედეგები.
მეექვსე თავში განხილულია შედეგების ინტერპრეტაცია, შეზღუდვები და ვალიდობის საფრთხეები.
მეშვიდე თავში ჩამოყალიბებულია დასკვნები და შემდგომი კვლევის მიმართულებები.

\chapter{ლიტერატურის მიმოხილვა}

\section{სამედიცინო გამოსახულების ანალიზის ტრადიციული მეთოდები}
სამედიცინო ვიზუალიზაციაში ტრადიციული ანალიზი დიდწილად ეყრდნობოდა წესებზე დაფუძნებულ და
ხელით მორგებულ მიდგომებს. პრაქტიკული მილსადენი, როგორც წესი, მოიცავდა:
(ა) წინასწარ დამუშავებას, (ბ) სეგმენტაციას, (გ) მახასიათებლების ამოღებას და (დ) კლასიფიკაციას.

\textbf{წინასწარი დამუშავება.}
ყველაზე გავრცელებული ოპერაციებია ხმაურის შემცირება (Gaussian/Median ფილტრები),
კონტრასტის ლოკალური გაძლიერება (ჰისტოგრამის ეკვალიზაცია, CLAHE), და ინტენსივობის ნორმალიზაცია.
ეს ნაბიჯები ხშირად აუმჯობესებს ვიზუალურ ხარისხს, თუმცა არ უზრუნველყოფს მოდელის მდგრადობას
დომენურ ცვლადობაზე.

\textbf{სეგმენტაცია.}
კლასიკური მეთოდები (ზღურბლოვანი ალგორითმები, კიდეების დეტექცია, region growing, watershed)
ეფექტურია მარტივ სცენარებში, მაგრამ მნიშვნელოვნად სუსტდება მაშინ, როცა გამოსახულება
ჰეტეროგენულია, საზღვრები ბუნდოვანია, ან არსებობს არტეფაქტები.

\textbf{ხელით ინჟინერირებული მახასიათებლები.}
ტრადიციულად, ტექსტურული, ფორმითი და ინტენსივობითი დესკრიპტორები გამოიყენებოდა
SVM/RF/ლოგისტიკურ კლასიფიკატორებთან ერთად.
ამ მიდგომების მთავარი შეზღუდვა არის მახასიათებლების ხელით შერჩევაზე მაღალი დამოკიდებულება,
რაც ამცირებს მასშტაბირებადობასა და განზოგადების უნარს.

\textbf{შეჯამება.}
ტრადიციულმა მეთოდებმა შექმნა ისტორიული საფუძველი CAD სისტემებისთვის,
მაგრამ რთულ მრავალლეიბლიან კლინიკურ ამოცანებში მათი ზღვარი აშკარაა.
ამიტომ, თანამედროვე კვლევა ფოკუსირებულია მონაცემებზე დაფუძნებულ ღრმა მოდელებზე.

\section{მანქანური სწავლება სამედიცინო გამოსახულებისთვის}
მანქანურმა სწავლამ სამედიცინო ანალიზი გადაიყვანა ფიქსირებული წესებიდან სტატისტიკურ სწავლაზე.
ძირითადი პარადიგმებია:

\begin{enumerate}[leftmargin=1.6em]
    \item \textbf{ზედამხედველობითი სწავლება}: ანოტირებულ მონაცემებზე კლასიფიკაცია/რეგრესია.
    \item \textbf{ზედამხედველობის გარეშე სწავლება}: კლასტერიზაცია, ანომალიების აღმოჩენა, განზომილების შემცირება.
    \item \textbf{ნახევრად ზედამხედველობითი სწავლება}: მცირე ლეიბლიან ნაკრებთან ერთად დიდი არალეიბლიანი მონაცემების გამოყენება.
\end{enumerate}

კლასიკურ ML მილსადენში წარმატება დამოკიდებულია feature engineering-ის ხარისხზე.
SVM, Random Forest, k-NN და ლოგისტიკური რეგრესია კვლავ სასარგებლო საბაზისო ხაზებია,
თუმცა რთული ვიზუალური პატერნების დაჭერა მათთვის შეზღუდულია.
ამ მიმართულებაში მნიშვნელოვანი თეორიული ჩარჩოები მოცემულია თანამედროვე სტატისტიკური სწავლების ლიტერატურაში
\footnote{მაგალითად, ზოგადი ML საფუძვლებისთვის იხ. \cite{hastie2009elements}.}.

\section{ღრმა სწავლება სამედიცინო გამოსახულებისთვის}
ღრმა სწავლების მთავარი უპირატესობაა \textit{end-to-end} სწავლა,
სადაც მოდელი პირდაპირ პიქსელებიდან სწავლობს იერარქიულ რეპრეზენტაციებს.

\subsection{CNN და მისი ევოლუცია}
AlexNet, VGG, Inception, ResNet და DenseNet ოჯახმა აჩვენა,
რომ სიღრმე, რეგულარიზაცია და არქიტექტურული ინოვაციები მნიშვნელოვნად ზრდის ვიზუალურ ხარისხს.
CNN-ები განსაკუთრებით ძლიერია ლოკალური ტექსტურული ნიშნების ამოცნობაში,
რაც რენტგენოგრაფიული ამოცანებისთვის კრიტიკულია.

\subsection{სეგმენტაციის არქიტექტურები}
FCN/U-Net კლასის მოდელებმა სემანტიკურ სეგმენტაციაში დე-ფაქტო სტანდარტი შექმნა.
U-Net-ის skip-კავშირები ეფექტურად აერთიანებს მაღალი გარჩევადობის ლოკალურ ინფორმაციას
და ღრმა ფენების სემანტიკურ კონტექსტს.

\subsection{ტრანსფერული სწავლება და Vision Transformer}
ტრანსფერული სწავლება აუცილებელია მაშინ, როცა სამედიცინო ანოტაციები შეზღუდულია.
ImageNet-ზე ან თვითზედამხედველურ რეჟიმში წინასწარ გაწვრთნილი backbone-ები მნიშვნელოვნად ამცირებს
სწავლის დროს და ზრდის სტაბილურობას.
Vision Transformer მიდგომები, განსაკუთრებით self-supervised ოჯახები (DINO/MAE),
ძლიერ შედეგებს აჩვენებს გლობალური კონტექსტის მოდელირებაში
\footnote{ViT-ის საფუძვლებისთვის იხ. \cite{dosovitskiy2020image}; self-supervised ოჯახისთვის იხ. \cite{caron2021dino,he2022mae,oquab2024dinofamily}.}.

\section{ChestMNIST-ზე მრავალნიშნიანი კლასიფიკაციის კვლევები}
გულმკერდის რენტგენოგრაფიის მრავალნიშნიანი კლასიფიკაცია წარმოადგენს პრაქტიკულად მნიშვნელოვან ამოცანას:
ერთი გამოსახულება შეიძლება შეიცავდეს ერთზე მეტ პათოლოგიას.
საჯარო ბენჩმარკები (ChestX-ray14, CheXpert, MIMIC-CXR, MedMNIST v2/ChestMNIST)
ქმნის შედარებად ექსპერიმენტულ ბაზას
\footnote{მონაცემთა ძირითადი წყაროები: \cite{wang2017chestxray14,irvin2019chexpert,johnson2019mimiccxr,yang2023medmnistv2}.}.

მიუხედავად ამისა, ამოცანა რთულდება შემდეგი ფაქტორებით:
კლასთა მკვეთრი დისბალანსი, ლეიბლების ხარისხის ცვალებადობა,
კო-მორბიდული პათოლოგიების თანაარსებობა და დომენური ცვლა სხვადასხვა კლინიკურ ცენტრს შორის.
ამ პირობებში მხოლოდ Accuracy არ არის საკმარისი;
აუცილებელია მაკრო-AUC/AUPRC, F1, კალიბრაცია და გაურკვევლობა.

\section{გამოწვევები და ღია საკითხები}
სამედიცინო AI-ის ტრანსლაციისას კრიტიკულია რამდენიმე ღია საკითხი:

\begin{enumerate}[leftmargin=1.6em]
    \item \textbf{მონაცემთა გამოწვევები}: დისბალანსი, ანოტაციის მაღალი ღირებულება, ჰეტეროგენული აკვიზიცია.
    \item \textbf{მოდელის გამოწვევები}: განზოგადება, მდგრადობა, ინტერპრეტირებადობა.
    \item \textbf{სანდოობის გამოწვევები}: overconfidence, calibration drift, uncertainty under shift.
    \item \textbf{კლინიკური ინტეგრაცია}: workflow-ში ჩაშენება, რეგულატორული შესაბამისობა, პასუხისმგებლობის განაწილება.
    \item \textbf{ეთიკა და მიკერძოება}: დემოგრაფიული/ინსტიტუციური bias და სამართლიანობის კონტროლი.
\end{enumerate}

კალიბრაციისა და uncertainty-ის ინტეგრირებული შეფასება განსაკუთრებით მნიშვნელოვანია,
რადგან მაღალი AUC ხშირად თანხვედრაში არ არის სანდო ალბათურ პროგნოზებთან
\footnote{კალიბრაცია და uncertainty: \cite{guo2017calibration,lakshminarayanan2017ensembles}; conformal გარანტიები: \cite{angelopoulos2021conformal}.}.

\section{კვლევითი ხარვეზი და ნაშრომის წვლილი}
ლიტერატურის ანალიზი აჩვენებს, რომ მრავალი კვლევა ფოკუსირდება მხოლოდ დისკრიმინაციულ მეტრიკებზე
(მაგ., AUC), მაგრამ ნაკლებად ფარავს სამ კომპონენტს ერთიანად:

\begin{enumerate}[leftmargin=1.6em]
    \item ონტოლოგიურად განპირობებული ლეიბლ-დამოკიდებულებები;
    \item long-tail ოპტიმიზაცია იშვიათი კლასების დასაცავად;
    \item სანდოობის სრული სტეკი (ensemble uncertainty + calibration + conformal coverage).
\end{enumerate}

წინამდებარე ნაშრომი სწორედ ამ სამ წერტილს აერთიანებს CXR-Synapse ჩარჩოში:
ELIXR-C/CXR-Foundation ვიზუალური რეპრეზენტაცია, RadLex-ზე აგებული ლეიბლ-გრაფი,
HLRF-ტიპის დანაკარგი, deep ensemble გაურკვევლობა, ECE/Brier კალიბრაცია და
Bonferroni-კორექტირებული conformal prediction.
რეპორტინგი მორგებულია TRIPOD/SPIRIT-AI პრინციპებზე
\footnote{რეპორტინგის ჩარჩოები: \cite{collins2015tripod,rivera2020spiritai}.}.

\section{თვითზედამხედველური სწავლება ვიზუალურ რეპრეზენტაციებში}
\textbf{Foundation representation learning.}
თვითზედამხედველური სწავლება განსაკუთრებით მნიშვნელოვანია სამედიცინო გამოსახულებებში,
რადგან მაღალი ხარისხის annotation ძვირი და დროით შეზღუდულია.
DINO/DINOv2 და MAE ტიპის მოდელებმა აჩვენეს, რომ დიდი რაოდენობის არალეიბლიანი გამოსახულებიდან
შესაძლებელია ისეთი feature space-ის მიღება, რომელიც downstream ამოცანებზე მცირე რაოდენობის label-ითაც მუშაობს.
რადიოლოგიაში ამ იდეას აგრძელებენ სპეციალიზებული CXR foundation models, მათ შორის ELIXR-C/CXR-Foundation,
რომლებიც უშუალოდ გულ-მკერდის გამოსახულებებზე მიღებულ representation-ს ქმნიან.

ამ ნაშრომში სწორედ ასეთი representation გამოიყენება: downstream მოდელი აღარ სწავლობს raw pixel-ებიდან,
არამედ იღებს 8x8 სივრცით feature map-ს. ეს ამცირებს გამოთვლით ხარჯს და საშუალებას იძლევა კვლევის ყურადღება
გადავიტანოთ label semantics-ზე, uncertainty-ზე და conformal safety-ზე.

\section{ტრანსფერული სწავლება სამედიცინო გამოსახულებებში}
ტრანსფერული სწავლების მთავარი არგუმენტი არის მონაცემთა ეკონომია.
სამედიცინო dataset-ები ხშირად მცირეა იმ მასშტაბთან შედარებით, რაც თანამედროვე ღრმა მოდელებს სჭირდება.
ImageNet-ზე გაწვრთნილი ResNet/DenseNet/EfficientNet მიდგომები წლების განმავლობაში წარმოადგენდა ძლიერ baseline-ს,
მაგრამ სამედიცინო domain shift ხშირად ზღუდავს მათ.
სამედიცინო foundation representation ამცირებს ამ პრობლემას, რადგან feature extractor უკვე სწავლობს
რადიოლოგიურ ინტენსივობებს, პოზიციურ კანონზომიერებებს და CXR-specific ვიზუალურ ტექსტურას.

CXR-Synapse არ აკეთებს ELIXR-C backbone-ის end-to-end fine-tuning-ს.
ეს არის გამიზნული არჩევანი: frozen feature რეჟიმი ზრდის რეპროდუცირებადობას,
ამცირებს hardware მოთხოვნებს და გამოკვეთს ამ ნაშრომის ძირითად წვლილს downstream classifier-ის დონეზე.

\section{ღრმა ensemble-ები და გაურკვევლობის რაოდენობრივი შეფასება}
ერთი deterministic მოდელი ხშირად ზედმეტად თავდაჯერებულია.
Deep ensemble ამცირებს ამ პრობლემას რამდენიმე დამოუკიდებელი seed-ით გაწვრთნილი მოდელის საშუალებით.
თუ მოდელები ერთსა და იმავე მაგალითზე განსხვავებულ probability score-ს აძლევენ, ეს epistemic uncertainty-ის
სიგნალია და შეიძლება გამოყენებულ იქნეს დამატებითი კლინიკური ყურადღების მისაქცევად.

მიმდინარე კოდში ensemble შედგება სამი seed-ისგან, ხოლო inference ეტაპზე dropout ფენები გააქტიურებულია
10 MC-pass-ისთვის. საბოლოო probability არის ყველა pass-ისა და ყველა seed-ის საშუალო.
ეს არ იძლევა სრულ Bayesian posterior-ს, მაგრამ პრაქტიკული და გამოთვლითად მისაღები uncertainty estimate-ია.

\section{კალიბრაცია და conformal პროგნოზი}
კალიბრაცია ნიშნავს, რომ პროგნოზირებული probability შეესაბამება რეალურ ემპირიულ სიხშირეს.
თუ მოდელი პათოლოგიას 0.8 ალბათობით აფასებს, მსგავსი შემთხვევების დაახლოებით 80\% უნდა იყოს დადებითი.
Expected Calibration Error (ECE) ამ განსხვავებას bin-ების მიხედვით ითვლის.

Conformal prediction ამატებს distribution-free უსაფრთხოების ფენას.
მრავალნიშნიან ამოცანაში იგი არ აბრუნებს მხოლოდ ერთ label-ს; იგი ქმნის prediction set-ს.
თუ set დიდია, ეს ნიშნავს, რომ მოდელი ფრთხილობს და რამდენიმე პათოლოგიას ტოვებს განხილვაში.
CXR-Synapse-ის micro-marginal conformal wrapper სწორედ ამ პრინციპით მუშაობს და validation set-ზე
არჩევს global multiplier-ს მიზნობრივი coverage-ის მისაღწევად.

\section{მრავალნიშნიანი სწავლება და კლასთა დისბალანსი}
CXR მონაცემებში ერთი გამოსახულება ხშირად შეიცავს რამდენიმე პათოლოგიას.
ამიტომ multiclass მიდგომა არ არის საკმარისი; საჭიროა multilabel learning, სადაც თითოეული label
დამოუკიდებელი binary target-ია, მაგრამ მათი ურთიერთკავშირი მაინც უნდა იქნას გათვალისწინებული.

კლასთა დისბალანსი განსაკუთრებით მწვავეა იშვიათ პათოლოგიებზე.
CXR-Synapse ამ პრობლემას პასუხობს ორი გზით: logit prior adjustment train prevalence-ის მიხედვით
და threshold optimization validation set-ზე. გარდა ამისა, graph-guided PaQ attention საშუალებას აძლევს
მოდელს იშვიათი label-ებიც დააკავშიროს სემანტიკურად ან სტატისტიკურად ახლო პათოლოგიებთან.

\iffalse
თვითზედამხედველური სწავლება (self-supervised learning, SSL) წარმოადგენს თანამედროვე პიონერულ მიდგომას,
რომელმაც და უშუალო აქტივაცია რეალიზებული იყოს წლების ხელით ანოტირებული ეტიკეტების გარეშე.
SSL განსაკუთრებით მნიშვნელოვანია სამედიცინო დომენში, სადაც ხარისხიანი ანოტაციები
დაკარგული და წვდება კოსტური.

Contrastive learning მიდგომები (SimCLR, MoCo) აკონცენტრირებენ
აკვიზიციის განმცვიფრებული გამოსახულებების შედარებაზე,
სადაც დაძლევითი მიზეზით განსხვავებული pairs მიმზიდველი და repulsive losses-ის საშუალებით.
DINO და უახლეს DINOv2 მოდელები აჩვენებენ განმეორებიან შედეგებს ImageNet-ზე,
რომელიც რადიოლოგიური დაავადებებისთვის განსაკუთრებით მნიშვნელოვანია,
რადგან ისინი აგებული არის მხოლოდ ბუნებრივი გამოსახულებების წინა განმწავლობაზე,
და ამიტომ გაშეყო ღრმა რეპრეზენტაციებს რომელიც ხელმოვარდილი ხდება სამედიცინო ტრანსფიკით.

Masked Autoencoders (MAE) მოდელი აგებულია დიდი ნაწილი გამოსახულების მოვილი პიქსელების დაფარვაზე,
შემდეგ მოდელი თავისთავადი აღმოცნობის დაკარგული მოტივაციის უბნების ხელმიჭ.
ეს approach უზრუნველყოფს განსაკუთრებით ძლიერი ტრანსფერული შედეგებს,
რადგან მოდელი სწავლობს მკრთალი რეპრეზენტაციებს მერე დილის ავტენთიკიტეტის ხელმიჭ.
MAE ფაქტორი შკოტსს ეკონომამდე (2-3x ფოა enim გამოთვლითი ცა ბაზედ SimCLR ან MoCo-თან)
მისი სკალებისთან უფრო დიდი მოდელების წინა განმწავლობა.

ეს ორი SSL მიდგომა (DINO/DINOv2 და MAE) განელი მკაფიოდ ელაგ ხელმეზობელი
დიდი რაოდენობის არალეიბლიანი რადიოლოგიური მონაცემების საფეხураზე.
სამედიცინო დომენებში, ეს method-ები ხშირად აჩვენებენ გაეფამის თანაბარი ან უმჯობესი ტრანსფერული შედეგი
გრუდიენტ-ორიენტირებული fine-tuning-ის საშუალებით გამოიცვლებიან პროტოკოლებით.

\section{ტრანსფერული სწავლება სამედიცინო გამოსახულებებში}
ტრანსფერული სწავლება რემის საფუძველი ხელმიჭ სამედიცინო AI-ზე,
მ克פაქტორი რადგან მოდელების მოთხოვნილი რაოდენობის მოეხელი
დაშეზღუდა ხელმოვარდილი ზომის სამედიცინო ბენჩმარკით (ხშირად 100K-500K).

ტრადიციული მიდგომა კრუმის ImageNet-ზე განწინიერებული backbone (ResNet-50, DenseNet-121, EfficientNet)
ის fine-tune რაოდენობის სამედიცინო ამოცანაზე, მაგ. ChestMNIST.
თუმცა, უფრო ეფექტურია SSL წინა განმწავლობის გამოყენება DINO ან MAE-სთან,
სადაც მოდელი უკვე იზი წინა განმწავლობა ზოგადი რადიოლოგიური და მედიცინური ფიზიკის თეორიით.

ფინ-ტიუნინგის ის მიდგომა თეოხელი დამოკიდებული არის
რა-რა ფაქტორების გაწინააღმდეგო აღქმის:

\begin{enumerate}[leftmargin=1.6em]
    \item \textbf{Layer-wise learning rate decay (LLRD)}: უფრო უმცროსი learning rate უფром ღრმა ფენებისთვის;
    \item \textbf{Early stopping}: ვალიდაციის მეტრიკის მივი სამედიცინო რაოდენობით გაჩერება;
    \item \textbf{Mixup და CutMix}: რეგულარიზაციის მოტივაციის ტექნიკა overfitting rescue-ის რადგინმდაბლა;
    \item \textbf{Stochastic Weight Averaging (SWA)}: საშუალო წონიკა აჭირობის რთული optimum-დან.
\end{enumerate}

უდიდეს მხარდამჭერ ან გადამკვეთი უდიდეს მოქმედებამდე
დაბალი learning rate ღრმა ფენებში დამოკიდებული აღქმა სან-ფოა ტიპის,
სადაც უბდეს რამდენიმე ფენი უფრო მუშაობაზე (საჭირო აკტივაცია ხელმიჭ ფიზიკაზე)
მხოლოდ რწელი feature მისაბამობის შემაარი უმცროსი დოზირებული ბელასი შესაძლო.
ეს ოჰო აღქმის ფაქტორი ხელი ერთიანი ხელმიდამ უდიდეს ფიზიკაში
უმცროსი degeneration უსაფრთხოების და უფრო არტისტული მიმახევების ხელმეზობელი.

\section{ღრმა ensemble-ები და გაურკვევლობის რაოდენობრივი}
პირიანი deterministic მოდელი ხშირად overconfident ხდება,
განსაკუთრებით კლასებში ღმკალი ხტომის სიხშირე დაბალი.
ღრმა ensemble მეთოდოლოგია ხელმოვარდილი რამდენიმე დამოუკიდებელი მოდელით,
სადაც თითოეული გაწვრთნილი მხოლოდ რეიტინგის ნიმუშების სოტირებული,
განსხვავებული random seed-ებით და პარამეტრების ინიციალიზაციით, რაც უფრო დიდი დისვარსია ხელმეზობელი.

ეპისტემური გაურკვევლობა (მოდელის უცნობი ავტობიოგრაფიით დაკავშირებული)
წარმოადგენილი როგორც ensemble-ის წევრებს შორის disagreement variance.
Aleatoric გაურკვევლობა (მონაცემთა თავისთავადი აბერაცია) 
ხელმოვარდილი როგორც predict entropy თითოეული enrollment-ის.
გაერთიანებული uncertainty შედეგი ხელმიჭ აძილებელი
უმჯობესი posterior ესტიმატი ensemble წევრებიდან.

მიუხედავად გამოთვლითი დანახარჯი (গ ensemble მოდელი gödz უსაფრთხო გამოქვაბითი დანახარჯი),
deep ensemble უპირატესია კლინიკურ აპლიკაციებში,
რადგან uncertainty რეიტინგი რომელიცო ხელი კლინიცისტებს კრიტიკულ მედიცინურ გადაწყვეტილებებში.
უდიდეს მხარდამჭერ კეთილი predicting uncertainty დამოვიტარებული
საბაზო რაოდენობის აბსტრაქტული मনพอสม-დ დაბალი-რისკ პაციენტებთან მაღალი კუთხე-მოძებნილებთან.

\section{კალიბრაცია და conformal პროგნოზი}
კალიბრაცია მოწყობილა უთითებს რამდენი კარგი მოდელის predicted პროვაბილიტი 
რეალურ empirical ფ predic არის დაკომიტირებული ხელმეზობელი ფიზიკაში.
Deep neural networks ხშირად overconfident ხდებიან საბაზო trained-დან,
განსაკუთრებით multilabel წრთილებში სადაც ئ სიტუაციული რელატური კერძოდ რჩება.

Expected Calibration Error (ECE) სტანდარტული მეტრიკა calibration რაოდენობრივი მოვიკიდეთ ხელმეზობელი აბსტრაქტული.
არსებული მეტრიკა ფიქსირებული დიაპაზონი ერთეულებლი და ხელმიდამ უსაფრთხოების მეტრიკა.
Temperature scaling არის მარტივი კალიბრაციის მეთოდი რომელიცო გამოიყენება post hoc უშუო რძეთ
مسაุ probabilistic გამოცემებში უდიდეს მხარდამჭერ deterministic პროცესი პირიანი ქვესფეროებში.

Conformal prediction უზრუნველყოფს distribution-free გარანტიებს prediction sets-ზე,
რომელმა აქვთ მჯერელი სტატისტიკური დაზღვევა თან კონკრეტული ნიჯო დაშვებებით.
Multilabel მხედელობაზე Bonferroni კორექცია გამოიყენება conservative approach უსაფრთხოების მდგრადობისთვის,
სადაც თითოეული მოქმედების კვანტილი არის ითვლება დამოუკიდებელი.
ეს მონაკვეთი უზრუნველყოფს რომ prediction sets დააკმაყოფილებს მნიშვნელოვანი coverage არეგულირებელი
უსაფრთხოების სეგმენტების ხელომიჭვით და ქვეღაირის გაშესუ.

\section{მრავალნიშნიანი სწავლება და კლასთა დისბალანსი}
რეალურ რადიოლოგიურ აპლიკაციებში, ერთი მედიცინური გამოსახულება ხშირად შეიცავს ერთზე მეტი პათოლოგია.
Multilabel კლასიფიკაციის ამოცანა ნიშნავს რომ თითოეული კლასი უნდა რეგულირდეს დამოუკიდებელი,
და ერთი ნიმუშს შეუძლია უნდა ჰქონდეს ერთზე მეტი positive labels ხელმოვარდილი აბსტრაქტული.

კლასთა დისბალანსი უდიდესი გამოწვევა multilabel პროცესებში:
ChestMNIST ბენჩმარკში pathology frequency 1\%-დან მილიონი 70\%-მდე,
მნიშვნელოვანი ღმკალი რაოდენობის თანხვედრაში უნივერსიტეტის მრავალი პაციენტი.
კლასთა დიპასტე დაწყებული რამდენიმე კლასის ზღვარი აღქმის ხელმეზობელი probabilistic ბალანსს:
ბელი ხელმოვარდილი ხელმიჭ განსაკუთრებული იმიტაციის loss ფუნქციად.

Asymmetric Loss და Focal Loss უგრძელი უდიდესი प্রमाणিত მეთოდებია long-tail გამოწვევებისთვის.
Asymmetric Loss უფრო დიდი პენალტი უხელსო დაწყებული false positives:
ხელი false positives ეტალონი ხელმეზობელი უსაფრთხოება ბელი ხელი false negatives უმეუნდა უდიდესი დაზღვევა ხელმიჭ.
Focal Loss მიდგომა პირიანი hard negatives რომელიცო უდიდეს ადამიანი არის დაკომიტირებული განსაკუთრებული უხელსო.
ორივე მიდგომი აჩვენებს მნიშვნელოვანი რაოდენობის გაუმჯობესება უდიდეს مرتبط imbalanced მონაცემებზე
აბსტრაქტული BCE loss-თან შედარებით, და ხშირად აერთიანებული optimized threshold-ებით
რათა კიდევ მეტი გაუმჯობესება რაოდენობის უფხელი-fone შესაძლო კლასებში.

\fi

\chapter{მეთოდოლოგია}

\section{მონაცემთა ბაზა და გაყოფის სტრატეგია}
კვლევის ძირითადი წყაროა MedMNIST v2-ის ChestMNIST მონაცემთა ბაზა,
რომელიც გულ-მკერდის რენტგენოგრამების 14-ლეიბლიან მრავალნიშნიან კლასიფიკაციას მოიცავს.
მიმდინარე კოდში მონაცემები იტვირთება \texttt{medmnist.ChestMNIST} კლასით, 224x224 ზომით,
ხოლო სასწავლო, ვალიდაციისა და სატესტო ნაწილები ინარჩუნებს MedMNIST-ის ოფიციალურ split-ებს.

მონაცემთა გაყოფის სანდოობისთვის გამოიყენება leakage-ის სასამართლო აუდიტი.
\texttt{data.py} თითოეული გამოსახულებისთვის ითვლის MD5 hash-ს და ამოწმებს,
რომ ვალიდაციისა და ტესტის ჩანაწერები არ მეორდებოდეს სასწავლო ნაწილში.
\texttt{audit.py} დამატებით მოიცავს cross-split hash-შემოწმებას, label co-occurrence ანალიზს,
სიგნალ-ხმაურის განაწილებას, label cardinality-ს და კლინიკური თანხვედრის საბაზისო წესებს.

\section{მონაცემთა წინასწარი დამუშავება}
მიმდინარე სისტემა end-to-end პიქსელურ ტრენინგს არ ასრულებს.
პირველ ეტაპზე \texttt{Extraction/extraction\_pipeline.py} ChestMNIST გამოსახულებებს გარდაქმნის
ELIXR-C/CXR-Foundation მოდელისთვის საჭირო ფორმატში.
გამოსახულება ინახება 16-bit PNG representation-ად, შემდეგ TensorFlow SavedModel-ის
\texttt{serving\_default} ხელმოწერით მიიღება სივრცითი feature map.

ექსპერიმენტული pipeline მუშაობს 10\%-იან decile subset-ზე (\texttt{FRACTION=0.10}),
რაც გამოთვლითი რესურსების შეზღუდვის პირობებში საშუალებას იძლევა შენარჩუნდეს სრული კვლევითი ციკლი:
ექსტრაქცია, სწავლება, ensemble შეფასება, bootstrap ტესტირება და conformal კალიბრაცია.
ფინალური ნიშნები ინახება \texttt{train\_embeddings.pt}, \texttt{val\_embeddings.pt} და
\texttt{test\_embeddings.pt} ფაილებში. თითოეულ ჩანაწერს აქვს 8x8 სივრცითი grid და 1376-განზომილებიანი
ELIXR-C feature vector, რაც \texttt{model.py}-ში აღიწერება \texttt{cxr\_dim=1376} პარამეტრით.

\section{მოდელის არქიტექტურა: CXR-Synapse}

CXR-Synapse აგებულია frozen-foundation feature-ებზე და შედგება ოთხი ძირითადი ნაწილისგან:
განზომილების შემცირება, 2D ანატომიური პოზიციური კოდირება, graph-guided Pathology-as-Query attention
და ბეტა-ევიდენციური პროგნოზირება.

\subsection{1. განზომილების შემცირება}
ELIXR-C-ის 1376-განზომილებიანი ნიშნები გადადის 384-განზომილებიან სამუშაო სივრცეში:
\begin{equation}
h = \text{LayerNorm}(\text{GELU}(W_2 \cdot \text{Dropout}(\text{GELU}(\text{LayerNorm}(W_1x))))).
\end{equation}
ეს ფენა ამცირებს პარამეტრების რაოდენობას და ამზადებს ნიშნებს pathology-query attention-ისთვის.

\subsection{2. 2D ანატომიური პოზიციური კოდირება}
მოდელი არ იყენებს ზოგად 1D positional embedding-ს.
\texttt{get\_2d\_sincos\_pos\_embed} ქმნის deterministic sine-cosine 2D კოდირებას 8x8 grid-ზე,
რაც ინარჩუნებს მარცხენა/მარჯვენა ფილტვისა და ზედა/ქვედა ანატომიური მდებარეობის ინფორმაციას.

\subsection{3. RadLex/BioViL-T სემანტიკური embedding-ები}
14 პათოლოგიის ტექსტური აღწერა წარმოდგენილია RadLex-ის შესაბამისი ტერმინებით.
\texttt{ensure\_radlex\_embeddings} იყენებს \texttt{microsoft/BiomedVLP-BioViL-T} მოდელს
768-განზომილებიანი embedding-ების მისაღებად. მიღებული embedding-ები ქმნის disease query-ების
საწყის სემანტიკურ ბაზას და ინახება \texttt{radlex\_embeddings\_14.pth} ფაილში.

\iffalse
\subsection{3. GraphGPS Classifier (ონტოლოგიურ განპირობებული)}
ლეიბლი განიხილება არა თვითმხარდამჭერი ალტერნატივებად, არამედ გა संযুक్తი სტრუქტურაში:

\begin{itemize}
    \item \textbf{Ontology Conditioning}: თითოეული 14 პათოლოგიის კვანძი $(v_k)$ თავის BioBERT RadLex ემბედინგით $(e_k)$ აღიჭურვება. ეს უზრუნველყოფს მომითამე კლინიკური სემანტიკის მოდელირებას;
    \item \textbf{Parallel მოთხოვნის Passing}: ორი პარალელური გზა:
    \begin{enumerate}
        \item \textbf{Local MPNN}: გრაფის კვანძებზე მუშაობა, სადაც ზღვარი აკავშირებს ხშირად თანაარსებოველ პათოლოგიებს (co-occurrence adjacency);
        \item \textbf{Global Transformer}: სახელმწიფო-level ატენშენი ნებისმიერი კვანძის ნებისმიერი სხვას კვანძამდე;
    \end{enumerate}
    \item \textbf{inal Logits}: თითოეული კლასის $k$ გამჯამებულ წარმოდგენა $W_k$ მოწოდება აღმდგენი მატრიცად, სადაც $(Q, K, V) = (z, z_{\text{cls}}, \mathcal{P})$, შემდეგ ლოგიტი დგება $\ell_k = z_k^\top W_k$.
\end{itemize}

\subsection{4. Student-t მდელიანი Space Clustering}
მშრალი კლასების ნიშნების სეგრეგაციაა ჩართული:
\begin{itemize}
    \item \textbf{Heavy-tailed იყო ცენტრი}: Student-t ის გამოყენება მაქსიმალური ხელმოკიდეობის დროს
    (DoF $= 4.0$) მდელიანი ფენის feature-space-იდან;
    \item \textbf{ატრიპლეტი Loss}: Batch-Hard ტრიპლეტი სამიზნის ანქორ დაღელვა განმუქცებული ფუძეები
    დამოწმების მობილური დომენიდან გარეთან.
\end{itemize}

\subsection{5. კალიბრირებული ინფერენსი: Deep Ensemble + TTA}
მხოლოდ stochastic კაპსულების (MC Dropout) ხარვეზის გაკვეთილმა:
\begin{itemize}
    \item $M=3$ დამოუკიდებელი მოდელი, თითოეული კოორდინირებული სიმი;
    \item Horizontal-flip Test-Time Augmentation (TTA);
    \item აქლიმატური კოვალნი მიერ ensemble variance, უზრუნველყოფილი epistemic თავისუფალი უნციტი.
\end{itemize}

\subsection{მრავალნიშნიანი head და ზღურბლების ტიუნინგი}
კლას-დამოუკიდებელი sigmoid ალბათობების შემდეგ,
თითოეული კლასის გადაწყვეტილების ზღურბლი ირჩევა ვალიდაციაზე
მაკრო-F1-ის მაქსიმიზაციით, ნაცვლად ფიქსირებული 0.5-ისა.

\fi

\subsection{4. Graph-guided Pathology-as-Query attention}
label dependency არ განიხილება დამოუკიდებელი sigmoid-ების უბრალო ნაკრებად.
\texttt{select\_adjacency\_threshold} co-occurrence მატრიცისთვის elbow-method-ით არჩევს ზღურბლს,
\texttt{build\_cooccurrence\_adjacency} კი ქმნის ნორმალიზებულ adjacency მატრიცას.
შემდეგ \texttt{PathologyCrossAttention} ასრულებს სამ ნაბიჯს:
\begin{enumerate}[leftmargin=1.6em]
    \item RadLex query-ების projection 384-განზომილებიან სივრცეში;
    \item graph message passing adjacency მატრიცით, რათა ხშირი თანაარსებობის მქონე პათოლოგიებმა ერთმანეთზე გავლენა მოახდინონ;
    \item cross-attention, სადაც query არის პათოლოგია, ხოლო key/value არის 8x8 ვიზუალური patch-ები.
\end{enumerate}
საბოლოო logit მიიღება normalized disease feature-სა და normalized query-ს შორის cosine similarity-ით,
რომელსაც ემატება learnable \texttt{logit\_scale}.

\subsection{5. Long-tail logit prior}
ChestMNIST-ში კლასთა დისბალანსი მკვეთრია.
\texttt{compute\_logit\_adjustment} თითოეული კლასისთვის ითვლის train prevalence-ის log-prior-ს:
\begin{equation}
\Delta_k = \tau \log \frac{\pi_k}{1-\pi_k}.
\end{equation}
მოდელში ეს prior ემატება ვიზუალურ logit-ს და ამცირებს ხშირი კლასების დომინირებას იშვიათ პათოლოგიებზე.

\section{სწავლების პროცესი და დანაკარგის ფუნქცია}
სწავლება ხორციელდება \texttt{train.py}-ში აღწერილი კონფიგურაციით:
35 epoch, batch size 64, AdamW optimizer, საწყისი learning rate $2 \times 10^{-4}$,
weight decay 0.05, 200-step warmup, cosine schedule, gradient clipping 1.0 და early stopping
ვალიდაციის AUC-ზე patience=10 პირობით.

ძირითადი დანაკარგია \texttt{StrictlyProperBetaEvidentialLoss}.
იგი აერთიანებს binary cross-entropy-ს და ბეტა-ევიდენციურ KL regularization-ს:
\begin{equation}
\mathcal{L} = \mathcal{L}_{BCE} + 0.05 \cdot \lambda_t \cdot \mathcal{L}_{KL},
\end{equation}
სადაც $\lambda_t$ იზრდება annealing წესით პირველი 20 epoch-ის განმავლობაში.
ლოგიტები გარდაიქმნება ბეტა განაწილების პარამეტრებად:
\begin{equation}
\alpha = \exp(z) + 1,\quad \beta = \exp(-z) + 1.
\end{equation}
ამით მოდელი სწავლობს არა მხოლოდ ალბათობას, არამედ პროგნოზის ევიდენციურ სიძლიერესაც.

LISA (Language-Invariant Semantic Anchoring) დამატებით ამაგრებს label semantics-ს:
საწყისი RadLex embedding-ების cosine topology გამოიყენება anchor-ად, ხოლო მიმდინარე query topology-ს
MSE დაშორება ემატება საერთო objective-ს წონით 0.5. ეს იცავს პათოლოგიური query-ების კლინიკურ
სემანტიკას გადაჭარბებული ტრენინგული drift-ისგან.

\iffalse

კვლევაში გამოიყენება ორფაზა სწავლების პროტოკოლი:

\subsection{Phase 1: Warm-up (3 ეპოქა)}
\begin{itemize}
    \item Backbone დაფიქსირებული;
    \item GraphGPS და Student-t Clustering heads საწყის მდგომარეობაში;
    \item დაბალი learning rate, მხოლოდ წინა თოკენი მოძრაობა.
\end{itemize}

\subsection{Phase 2: Fine-tuning (12 ეპოქა)}
\begin{itemize}
    \item ფult end-to-end ბეკპროპაგაცია;
    \item Layer-Wise Learning Rate Decay (LLRD): decay factor = 0.75, base LR = $5 \times 10^{-5}$;
    \item ეს დაზოგავს DINOv3-ის მაღალი დონის სივრცითი პრიორებს.
\end{itemize}

\subsection{H-FZLPR Loss (Hierarchical Focal Zero-Reference Pairwise Ranking)}

სამ კომპონენტი აერთიანებულია:

\begin{enumerate}[leftmargin=1.6em]
    \item \textbf{Focal Weighted Cross-Entropy}: დაკმაყოფილების მოდულაციის ფოკუსირებული რეგიონებზე:
    $$\mathcal{L}_{\text{focal}} = -\alpha(1-p_t)^{\gamma} \log(p_t),$$
    სადაც $\gamma=2.0$ (ფოკუსირების უფლებამოსილების ხელმოწერა), 
    და $\alpha$ გაბნიკი ურთული რეტელებისთვის;
    
    \item \textbf{Pairwise Ranking (ZLPR)}: logsumexp ფუნქცია, რომელიც რეგულარიზებს რელატიური ლოგიტების რანჟირებას:
    $$\mathcal{L}_{\text{pair}} = \log \sum_{i : y_i = 0} \exp(\ell_i - \ell_j \text{ for } y_j = 1);$$
    
    \item \textbf{Hierarchical Penalty}: კლინიკური იერარქიის გამოკვლევა (მაგ., Pneumothorax უნდა ჩაითვალოს მედიცინესთან) პენალტიზაციით logical violations:
    $$\mathcal{L}_{\text{hier}} = \sum_k \text{ReLU}(\ell_k - \ell_j) \quad \text{(აკრძალული ლოგიცილი სხვობა)}.$$
\end{enumerate}

საერთო objective:
$$\mathcal{L}_{\text{total}} = \mathcal{L}_{\text{focal}} + \lambda_h \mathcal{L}_{\text{hier}} + \lambda_t \mathcal{L}_{\text{triplet}},$$
სადაც $\lambda_h = 0.1$, $\lambda_t = 0.05$.

\subsection{Optimizer და Scheduler}

\begin{itemize}
    \item \textbf{AdamW}: $\beta_1 = 0.9$, $\beta_2 = 0.999$, weight decay = 0.05;
    \item \textbf{Learning Rate Schedule}: linear warmup (500 steps) + cosine annealing;
    \item \textbf{Stochastic Weight Averaging (SWA)}: დაწყება 10-ე ეპოქიდან, გაუმჯობესებული generalization და ალბათური კალიბრაცია;
    \item \textbf{Early Stopping}: ვალიდაციის მაკრო-F1 ზე, patience=5 ეპოქა.
\end{itemize}

\fi

\section{SWA, ensemble, კალიბრაცია და conformal პროგნოზი}
25-ე epoch-იდან გამოიყენება Stochastic Weight Averaging (SWA), \texttt{swa\_lr=5e-5} პარამეტრით.
საბოლოო შეფასება ეფუძნება სამ დამოუკიდებელ seed-ს: 42, 43 და 44.
\texttt{DeepEnsembleTTAEvaluator} თითოეული მოდელისთვის inference-ისას ააქტიურებს dropout ფენებს
და ატარებს 10 MC-pass-ს. მიღებული პროგნოზები ჯერ საშუალოვდება მოდელის შიგნით,
შემდეგ კი ensemble წევრებს შორის. ეს იძლევა epistemic uncertainty-ის პრაქტიკულ შეფასებას.

კალიბრაცია ფასდება Expected Calibration Error-ით (ECE), 15 bin-იანი per-class სქემით.
გადაწყვეტილების ზღურბლები ირჩევა ვალიდაციის მონაცემებზე \texttt{optimise\_thresholds} ფუნქციით,
რომელიც თითოეული კლასისთვის probability percentile grid-ზე ეძებს საუკეთესო F1-ს.

უსაფრთხოების დამატებითი ფენა არის \texttt{MultiLabelConformalPredictor}.
იგი ვალიდაციის set-ზე არჩევს ერთ global multiplier-ს, რომელიც ამცირებს ან ზრდის კლასობრივ ზღურბლებს ისე,
რომ micro-marginal coverage მიაღწიოს მიზნობრივ $1-\alpha=0.90$ დონეს.
სატესტო ეტაპზე conformal set განისაზღვრება წესით:
\begin{equation}
S(x)=\{k: p_k(x) \geq m \cdot t_k\},
\end{equation}
სადაც $t_k$ არის ვალიდაციაზე ოპტიმიზებული კლასობრივი threshold, ხოლო $m$ conformal multiplier.

\iffalse
ეპისტემური გაურკვევლობისთვის გამოიყენება $M=5$ დამოუკიდებელი მოდელის deep ensemble.
ინფერენსისას ითვლება როგორც საშუალო პროგნოზი, ისე disagreement variance.

კალიბრაციის შეფასება ეფუძნება ECE და Brier მეტრიკებს.
distribution-free გარანტიებისთვის გამოიყენება conformal prediction,
სადაც calibration-set კვანტილი განსაზღვრავს prediction set-ის სიგანეს.
მრავალნიშნიან ამოცანაზე გამოიყენება Bonferroni კორექცია.

\fi

\section{შეფასების მეტრიკები და ექსპერიმენტის დიზაინი}
ძირითადი მეტრიკებია:
macro AUC, macro F1, mean per-class ECE, conformal coverage და prediction set size.

ექსპერიმენტის დიზაინი აერთიანებს:
\begin{enumerate}[leftmargin=1.6em]
    \item დეტერმინისტულ run-manifest-ს (seed, split, config, commit hash);
    \item 2000-ჯერად bootstrap ნდობით ინტერვალებს;
    \item დაწყვილებულ bootstrap მნიშვნელოვნების ტესტირებას ერთ-მოდელიან baseline-სა და CXR-Synapse ensemble-ს შორის;
    \item ვალიდაციაზე დაფიქსირებულ threshold-ებს, რათა ტესტზე არ მოხდეს ხელახალი მორგება.
\end{enumerate}

აღნიშნული მეთოდოლოგია უზრუნველყოფს შედეგების სანდოობას,
რეპროდუცირებადობას და კლინიკური ინტერპრეტაციისთვის საჭირო გამჭვირვალობას.

\chapter{Experimental Protocol and Reproducibility}
\section{Hardware and Software}
\begin{table}[H]
\centering
\caption{Recommended environment description.}
\begin{tabular}{p{0.34\textwidth}p{0.55\textwidth}}
\toprule
Item & Value \\
\midrule
GPU & NVIDIA CUDA GPU (exact model and VRAM captured in run manifest) \\
CPU & Multi-core x64 CPU (exact model logged in run manifest) \\
RAM & Experiment machine memory captured in run manifest \\
OS & Windows (64-bit) \\
Python & 3.13.12 (virtual environment) \\
Core libraries & PyTorch, torchvision, transformers, numpy, scipy, pandas (exact versions in manifest) \\
\bottomrule
\end{tabular}
\end{table}

\section{Determinism and Run Manifest}
Each run stores:
\begin{enumerate}[leftmargin=1.6em]
    \item global seed and ensemble seeds,
    \item exact split manifest,
    \item code commit hash,
    \item config file and hyperparameters,
    \item checkpoints and evaluation artifacts.
\end{enumerate}

\section{Evaluation Metrics}
Headline metrics:
macro AUC, macro AUPRC, macro/micro F1,
per-class AUC and AUPRC,
ECE, Brier score,
and conformal coverage-set size statistics.

\section{Statistical Inference}
\subsection{Bootstrap Confidence Intervals}
For each metric $m$, 95\% bootstrap CI is reported from repeated resampling.

\subsection{Paired Significance Testing}
Differences between baseline and CXR-Synapse are evaluated using paired bootstrap
with two-sided p-values and confidence intervals over metric deltas.

\section{Ablation and Comparator Matrix}
\begin{table}[H]
\centering
\caption{Minimum ablation matrix required for thesis-level claims.}
\begin{tabular}{p{0.07\textwidth}p{0.74\textwidth}p{0.14\textwidth}}
\toprule
ID & Configuration & Status \\
\midrule
A0 & Baseline encoder + BCE & Required comparator \\
A1 & A0 + ontology graph classifier & Required comparator \\
A2 & A1 + long-tail loss & Required comparator \\
A3 & A2 + threshold optimization & Required comparator \\
A4 & A3 + deep ensemble uncertainty & Required comparator \\
A5 & A4 + conformal prediction & Required comparator \\
\bottomrule
\end{tabular}
\end{table}

\chapter{შედეგები}
წინამდებარე თავში წარმოდგენილია \texttt{main.py}-ში აღწერილი მიმდინარე ექსპერიმენტული pipeline-ის
საბოლოო შეფასება. შედეგები ეფუძნება ChestMNIST-ის სატესტო ნაწილს, ხოლო threshold-ები,
conformal multiplier და calibration logic ფიქსირდება ვალიდაციის ნაწილზე.

\section{საბოლოო რაოდენობრივი შედეგები}
\begin{table}[H]
\centering
\caption{CXR-Synapse ensemble-ის მიმდინარე შედეგები ChestMNIST სატესტო ნაკრებზე}
\begin{tabular}{p{0.52\textwidth}p{0.32\textwidth}}
\toprule
მეტრიკა & მნიშვნელობა \\
\midrule
Macro AUC & 0.7913 \\
AUC 95\% CI & გენერირდება 2000-ჯერადი bootstrap-ით \\
$\Delta$AUC p-value & გენერირდება paired bootstrap test-ით \\
Macro F1 & გენერირდება ოპტიმიზებული threshold-ებით \\
Mean ECE & 0.0120 \\
Conformal coverage & 92.4\% \\
Average prediction set size & გენერირდება \texttt{main.py}-ის ბოლო summary-ში \\
\bottomrule
\end{tabular}
\label{tab:cxr_synapse_current_results}
\end{table}

Macro AUC 0.7913 მიუთითებს, რომ მოდელი ინარჩუნებს სტაბილურ დისკრიმინაციულ უნარს 14 პათოლოგიის
საშუალო ჭრილში. ამავდროულად, mean ECE 0.0120 განსაკუთრებით მნიშვნელოვანია,
რადგან სამედიცინო decision-support სისტემისთვის სწორი კალიბრაცია ხშირად ისეთივე კრიტიკულია,
როგორც მაღალი ranking performance. Conformal coverage 92.4\% აჩვენებს, რომ უსაფრთხოების wrapper
აღწევს მიზნობრივ 90\%-ზე მაღალ micro-marginal დაფარვას.

\section{შედარება ერთ-მოდელიან baseline-თან}
შედარების baseline არის ბოლო checkpoint-იდან ჩატვირთული ერთი \texttt{CXR\_Synapse\_Foundation} მოდელი,
რომელიც იგივე adjacency mask-ს, RadLex embedding-ებს და logit prior-ს იყენებს.
საბოლოო pipeline ამ baseline-ს ადარებს სამწევრიან deep ensemble-ს MC-Dropout inference-ით.
paired bootstrap test-ის მიზანია შეაფასოს, არის თუ არა ensemble პროგნოზის AUC-გაუმჯობესება
სტატისტიკურად სარწმუნო იმავე სატესტო მაგალითებზე.

\section{კალიბრაციის ანალიზი}
ECE ითვლება თითოეული კლასისთვის ცალ-ცალკე და შემდეგ საშუალოვდება.
დაბალი ECE ადასტურებს, რომ მოდელის ალბათობები პრაქტიკულად ახლოსაა ემპირიულ სიხშირეებთან.
ეს დასკვნა განსაკუთრებით მნიშვნელოვანია იმ გარემოში, სადაც საბოლოო მომხმარებელმა შეიძლება
გამოიყენოს მოდელის score არა როგორც დიაგნოზი, არამედ როგორც triage ან მეორე მოსაზრების სიგნალი.

\section{Conformal prediction}
\texttt{MultiLabelConformalPredictor} ქმნის prediction set-ს თითოეული გამოსახულებისთვის.
სისტემა არ აიძულებს მოდელს ერთ საბოლოო დიაგნოზზე, არამედ აბრუნებს იმ პათოლოგიების ნაკრებს,
რომელთა გამორიცხვაც მოცემული რისკის დონეზე არ არის უსაფრთხო.
მიღებული 92.4\%-იანი დაფარვა აჩვენებს, რომ მოდელი იშვიათი და გაურკვეველი შემთხვევებისთვის
უფრო ფრთხილ პროგნოზს იძლევა, რაც კლინიკურ გარემოში სასურველი თვისებაა.

\section{ხარისხობრივი ანალიზისთვის საჭირო მასალები}
საბოლოო დაცვის ვერსიაში რეკომენდებულია დაემატოს შემდეგი ფიგურები, რომლებიც უშუალოდ მიმდინარე კოდიდან
შეიძლება გენერირდეს:
\begin{enumerate}[leftmargin=1.6em]
    \item per-class ROC ან PR მრუდები;
    \item calibration reliability diagram და per-class ECE bar plot;
    \item conformal set size-ის განაწილება;
    \item forensic audit-ის heatmap-ები \texttt{audit.py}-დან.
\end{enumerate}
ამ ეტაპზე სეგმენტაციის Dice/IoU მეტრიკები არ გამოიყენება, რადგან პროექტი კლასიფიკაციის ამოცანაა
და segmentation head კოდში არ არსებობს.

\section{შეზღუდვები შედეგების ინტერპრეტაციაში}
შედეგები უნდა წაიკითხოს სამი პირობის გათვალისწინებით:
\begin{enumerate}[leftmargin=1.6em]
    \item ექსპერიმენტი მუშაობს ChestMNIST-ის 10\%-იან extracted subset-ზე;
    \item ELIXR-C feature extraction frozen რეჟიმშია და არა end-to-end fine-tuning;
    \item გარე კლინიკური ვალიდაცია CheXpert/MIMIC-CXR-ზე ჯერ არ არის შესრულებული მიმდინარე კოდში.
\end{enumerate}
შესაბამისად, ნაშრომის დასკვნა ეხება კვლევით პროტოტიპს და არა მზა კლინიკურ პროდუქტს.

\iffalse
წინამდებარე თავში წარმოდგენილია მე-2 თავში აღწერილი მეთოდოლოგიის საფუძველზე ჩატარებული
ექსპერიმენტების შედეგები. აქცენტი გაკეთებულია შემუშავებული მოდელის \textbf{MultiLabelModel}
(\textbf{EfficientNet-B0} ბექბონით) ეფექტურობის რაოდენობრივ და ხარისხობრივ შეფასებაზე,
ChestMNIST-ის ხელშეუხებელ სატესტო ნაკრებზე. წარმოდგენილი შედეგები ასახავს მოდელის
განზოგადების უნარს უნახავ მონაცემებზე.

\section{რაოდენობრივი შედეგები (Quantitative Results)}
მოდელის საბოლოო შეფასება განხორციელდა მხოლოდ სატესტო ნაკრებზე, 2.5 პუნქტში განსაზღვრული მეტრიკებით.
შეფასებისთვის გამოყენებულია ვალიდაციაზე არჩეული საუკეთესო checkpoint (ადრეული შეჩერების კრიტერიუმით).

\subsection{საერთო ეფექტურობის შეფასება}
ცხრილში \ref{tab:main_results} მოცემულია MultiLabelModel-ის ძირითადი შედეგები.

\begin{table}[H]
\centering
\caption{MultiLabelModel-ის (EfficientNet-B0) ეფექტურობა ChestMNIST სატესტო ნაკრებზე}
\begin{tabular}{p{0.62\textwidth}p{0.22\textwidth}}
\toprule
მეტრიკა & მნიშვნელობა \\
\midrule
\multicolumn{2}{l}{\textbf{მრავალნიშნიანი კლასიფიკაცია (14 კლასი)}} \\
AUC-ROC (macro-average) & 0.7955 \\
F1-ქულა (macro-average, ოპტიმიზებული ზღურბლებით) & 0.2413 \\
\midrule
\multicolumn{2}{l}{\textbf{ბინარული კლასიფიკაცია (ნებისმიერი პათოლოგია vs. არც ერთი)}} \\
AUC-ROC & 0.7418 \\
F1-ქულა (ოპტიმიზებული ზღურბლით) & 0.6950 \\
\bottomrule
\end{tabular}
\label{tab:main_results}
\end{table}

\noindent შედეგების ინტერპრეტაცია:
\begin{enumerate}[leftmargin=1.6em]
    \item \textbf{მრავალნიშნიანი AUC-ROC = 0.7955} მიუთითებს კარგ დისკრიმინაციულ უნარზე 14 კლასისთვის საშუალოდ.
    \item \textbf{მრავალნიშნიანი F1 = 0.2413} ასახავს კლასთა დისბალანსის მაღალ გავლენას და იშვიათი კლასების
    ამოცნობის სირთულეს.
    \item \textbf{ბინარული AUC-ROC = 0.7418} მიუთითებს, რომ მოდელი უკეთ არჩევს "პათოლოგიური" და "არაპათოლოგიური"
    შემთხვევების ზოგად განცალკევებას.
    \item \textbf{ბინარული F1 = 0.6950} მრავალნიშნიან F1-ზე მნიშვნელოვნად მაღალია, რაც ადასტურებს,
    რომ მოდელი უფრო ეფექტურია "ნებისმიერი პათოლოგიის" იდენტიფიკაციაში, ვიდრე ყველა კონკრეტული ლეიბლის
    თანაბრად ზუსტ პროგნოზში.
\end{enumerate}

\subsection{შედარებითი ანალიზი (Comparative Analysis)}
ამ კვლევის საბოლოო რეპორტინგი ეფუძნება ერთ ძირითად კონფიგურაციას (MultiLabelModel, EfficientNet-B0),
ამიტომ სხვადასხვა არქიტექტურას შორის ერთსა და იმავე პროტოკოლში სრულმასშტაბიანი შედარება
ამ თავში არ არის წარმოდგენილი. თუმცა შედარებითი ჭრილის სახით შესაძლებელია დავაფიქსიროთ,
რომ ბინარული ამოცანის მეტრიკები (AUC/F1) მაღლა დგას მრავალნიშნიან მაკრო-F1-ზე,
რაც მოსალოდნელია 14-კლასიანი დისბალანსირებული კლასიფიკაციის პირობებში.

\subsection{დეტალური კლასიფიკაციის ანალიზი}
\textbf{შეცდომათა მატრიცა.} დეტალური შეცდომების სტრუქტურისთვის გამოიყენება ნორმალიზებული confusion matrix.

\begin{figure}[H]
\centering
\fbox{\parbox{0.92\textwidth}{\centering
ჩასვით ნორმალიზებული confusion matrix (heatmap) სატესტო ნაკრისთვის.\\
რეკომენდებული ბილიკი: \texttt{docs/figures/ch3_confusion_matrix.png}
}}
\caption{ნორმალიზებული შეცდომათა მატრიცა სატესტო ნაკრებზე}
\label{fig:ch3_confusion}
\end{figure}

\textbf{ROC მრუდი.} მოდელის ზღურბლზე-დამოუკიდებელი შეფასებისთვის გამოიყენება ROC მრუდი.

\begin{figure}[H]
\centering
\fbox{\parbox{0.92\textwidth}{\centering
ჩასვით ROC მრუდი (multilabel macro ან binary).\\
რეკომენდებული ბილიკი: \texttt{docs/figures/ch3_roc_curve.png}
}}
\caption{ROC მრუდი სატესტო ნაკრებზე}
\label{fig:ch3_roc}
\end{figure}

\textbf{Precision-Recall მრუდი.} კლასთა დისბალანსის პირობებში PR მრუდი დამატებით ინფორმაციულია.

\begin{figure}[H]
\centering
\fbox{\parbox{0.92\textwidth}{\centering
ჩასვით Precision-Recall მრუდი (macro ან per-class).\\
რეკომენდებული ბილიკი: \texttt{docs/figures/ch3_pr_curve.png}
}}
\caption{Precision-Recall მრუდი სატესტო ნაკრებზე}
\label{fig:ch3_pr}
\end{figure}

\subsection{დეტალური სეგმენტაციის ანალიზი (არარელევანტური ამოცანა)}
წინამდებარე კვლევა ფოკუსირებულია \textbf{მრავალნიშნიან კლასიფიკაციაზე} და არა სეგმენტაციაზე.
ამიტომ Dice/IoU ტიპის სეგმენტაციის მეტრიკები არ გამოიყენება და ეს ქვეთავი წარმოდგენილია
მხოლოდ მეთოდოლოგიური სრული სტრუქტურის შესანარჩუნებლად.

\section{ხარისხობრივი შედეგები და ვიზუალიზაცია (Qualitative Results and Visualization)}
რაოდენობრივი მეტრიკების პარალელურად ხარისხობრივი ანალიზი მნიშვნელოვანია იმის გასაგებად,
თუ რა ტიპის შემთხვევებში არის მოდელი სტაბილური და სად ჩნდება შეცდომები.

\subsection{წარმატებული და წარუმატებელი პროგნოზების მაგალითები}
რეკომენდებულია სატესტო ნაკრებიდან წარმოდგენილი იყოს მინიმუმ ოთხი ტიპური შემთხვევა:
\textbf{True Positive}, \textbf{True Negative}, \textbf{False Positive}, \textbf{False Negative}.
თითოეულ მაგალითზე უნდა იყოს მითითებული:
(ა) რეალური ლეიბლი, (ბ) მოდელის პროგნოზი, (გ) შესაბამისი ალბათობა.

\begin{figure}[H]
\centering
\fbox{\parbox{0.92\textwidth}{\centering
ჩასვით 2x2 პანელი (TP/TN/FP/FN შემთხვევები) შესაბამისი წარწერებით.\\
რეკომენდებული ბილიკი: \texttt{docs/figures/ch3_case_examples.png}
}}
\caption{წარმატებული და შეცდომითი პროგნოზების მაგალითები}
\label{fig:ch3_cases}
\end{figure}

\subsection{ინტერპრეტაციის მეთოდების შედეგები}
მოდელის ახსნადობის გასაზრდელად მიზანშეწონილია Grad-CAM ან მსგავსი მეთოდის გამოყენება,
რათა დადგინდეს, რომელი ანატომიური რეგიონები ახდენენ გავლენას პროგნოზზე.

\begin{figure}[H]
\centering
\fbox{\parbox{0.92\textwidth}{\centering
ჩასვით Grad-CAM/attention heatmap მაგალითები (სწორი და შეცდომითი შემთხვევები).\\
რეკომენდებული ბილიკი: \texttt{docs/figures/ch3_gradcam_examples.png}
}}
\caption{ინტერპრეტაციის ვიზუალიზაცია (Grad-CAM ან ანალოგი)}
\label{fig:ch3_xai}
\end{figure}

ხარისხობრივი დაკვირვებები, როგორც წესი, აჩვენებს, რომ სწორი კლასიფიკაციისას მოდელის ყურადღება
კონცენტრირებულია კლინიკურად რელევანტურ რეგიონებზე, ხოლო შეცდომით პროგნოზებში ხშირად ფიქსირდება
არტეფაქტებზე ან დაბალკონტრასტულ უბნებზე გადატანილი ყურადღება.

\section{სტატისტიკური ანალიზი (Statistical Analysis)}
ეს ქვეთავი აფასებს წარმოდგენილი შედეგების სანდოობას და განმარტავს სტატისტიკური დასკვნის საზღვრებს.

\subsection{შედარების სტატისტიკური მნიშვნელოვნება}
რადგან საბოლოო რეპორტინგი გაკეთებულია ერთი ძირითადი მოდელით, არქიტექტურათშორისი წყვილური
ტესტები (მაგ., Wilcoxon/paired t-test) ამ თავში არ არის გამოყენებული.
სრული head-to-head შედარების შემთხვევაში ტესტირება უნდა ჩატარდეს იდენტურ split/protocol-ზე.

\subsection{სანდოობის ინტერვალები}
ამ თავში წარმოდგენილია სატესტო ნაკრებზე მიღებული \textit{წერტილოვანი} შეფასებები.
შემდგომი გამოცემისთვის რეკომენდებულია bootstrap-ზე დაფუძნებული 95\%-იანი ნდობითი ინტერვალების
დამატება AUC/F1 მეტრიკებისთვის, რათა შეფასდეს შედეგების ვარიაბელობა.

\subsection{ქვეჯგუფების ანალიზი}
ChestMNIST ფორმატში დემოგრაფიული ატრიბუტები სრულად ხელმისაწვდომი არ არის,
რის გამოც ასაკის/სქესის მიხედვით სანდო ქვეჯგუფური შეფასება ამ ვერსიაში შეზღუდულია.
მომდევნო ეტაპზე რეკომენდებულია გარე კოჰორტებზე სტრატიფიცირებული ანალიზის ჩატარება.

\subsection{სტატისტიკური ანალიზის შეზღუდვები}
\begin{enumerate}[leftmargin=1.6em]
    \item წარმოდგენილი მნიშვნელობები არის ერთი ფინალური run-ის შეფასება.
    \item მრავალნიშნიან კლასიფიკაციაში მაკრო-F1 განსაკუთრებით მგრძნობიარეა იშვიათ კლასებზე.
    \item ერთიანი პროტოკოლით მრავალმოდელური შედარების არქონა ზღუდავს მიზეზობრივ ინტერპრეტაციას.
\end{enumerate}

მიუხედავად აღნიშნული შეზღუდვებისა, მიღებული AUC/F1 შედეგები ადასტურებს,
რომ EfficientNet-B0-ზე აგებული MultiLabelModel წარმოადგენს ტექნიკურად გამართულ საბაზისო ჩარჩოს
გულმკერდის მრავალნიშნიანი ავტომატური ანალიზისთვის.

\subsection{Per-Class Performance Breakdown}
Detailed per-class AUC, AUPRC, and F1 scores are computed to assess which pathologies the model predicts well.
A per-class breakdown is particularly important because:
\begin{enumerate}[leftmargin=1.6em]
    \item Clinical significance varies by pathology (e.g., pneumothorax vs. scarring);
    \item Class imbalance is severe: rare pathologies may have low F1 despite reasonable AUC;
    \item F1 scores indicate whether recall on rare classes is acceptable for clinical deployment.
\end{enumerate}

Example per-class metrics:
\begin{itemize}
    \item High-frequency pathologies: AUC $\geq 0.80$, F1 $\geq 0.60$;
    \item Medium-frequency: AUC $0.70$-$0.80$, F1 $0.40$-$0.60$;
    \item Rare pathologies: AUC $0.60$-$0.75$, F1 $< 0.40$ (inherent to imbalance).
\end{itemize}

The aggregate macro-AUC (0.7955) indicates competitive discrimination, while macro-F1 (0.2413)
reflects the severe challenge of minority-class prediction.

\subsection{Ensemble Agreement and Epistemic Uncertainty}
Deep ensemble disagreement is examined:
\begin{enumerate}[leftmargin=1.6em]
    \item Average pairwise KL divergence among ensemble members: $\approx 0.08$ nats;
    \item Variance (epistemic uncertainty) correlates moderately with prediction error (Spearman $\rho \approx 0.35$);
    \item Samples with high disagreement: often near decision boundaries---clinically valuable for flagging borderline diagnoses.
\end{enumerate}

\subsection{Conformal Prediction Coverage}
Conformal prediction evaluated on test set with target coverage $\alpha = 0.95$:
\begin{enumerate}[leftmargin=1.6em]
    \item Empirical coverage: $95.2\%$ (close to target);
    \item Average set size: $5.3$ out of $14$ classes per sample;
    \item Interpretation: Small sets indicate high-confidence predictions; large sets indicate uncertainty or multi-pathology.
\end{enumerate}

\subsection{Reproducibility and Implementation Details}

\subsubsection{Software Environment}
Key dependencies:
\begin{enumerate}[leftmargin=1.6em]
    \item PyTorch 2.0+, torchvision 0.15+;
    \item MedMNIST v2 dataset;
    \item scikit-learn for metrics;
    \item matplotlib, seaborn, numpy, pandas.
\end{enumerate}

\subsubsection{Hardware and Runtime}
\begin{itemize}
    \item GPU: NVIDIA RTX 3090 (24GB VRAM);
    \item Training per ensemble member: ~8 hours;
    \item Total ensemble training: 40 hours;
    \item Inference per batch (128 samples): ~45 milliseconds.
\end{itemize}

\subsubsection{Data Augmentation}
Training set augmented with:
\begin{enumerate}[leftmargin=1.6em]
    \item Random horizontal flips (p=0.5);
    \item Random rotations ($\pm 10°$);
    \item Random brightness/contrast shifts;
    \item Mixup and CutMix.
\end{enumerate}

No augmentation during validation/test.

\fi

\chapter{დისკუსია}
მიღებული შედეგები აჩვენებს, რომ CXR-Synapse-ის მთავარი ღირებულება მხოლოდ AUC-ის გაუმჯობესებაში არ არის.
სისტემა აერთიანებს ცოდნაზე დაფუძნებულ label-გრაფს, RadLex/BioViL-T სემანტიკას, 2D სივრცით კოდირებას,
ევიდენციურ loss-ს და conformal safety wrapper-ს. ასეთი კომბინაცია პასუხობს სამედიცინო AI-ის
სამ ძირითად მოთხოვნას: დისკრიმინაცია, კალიბრაცია და გაურკვევლობის მართვა.

\section{შედეგების ინტერპრეტაცია}
Macro AUC 0.7913 ChestMNIST-ის 14-ლეიბლიან ამოცანაზე მიუთითებს კონკურენტულ ranking quality-ზე,
მაგრამ ნაშრომის პრაქტიკული წვლილი განსაკუთრებით ჩანს ECE=0.0120 მაჩვენებელში.
სწორად კალიბრირებული ალბათობა კლინიკურ workflow-ში უფრო სანდო triage-ს ქმნის,
რადგან მაღალი score ნაკლებად არის ცრუ თავდაჯერებულობის შედეგი.

Conformal coverage 92.4\% აჩვენებს, რომ მოდელი გაურკვეველ შემთხვევებში პროგნოზის ნაკრებს აფართოებს.
ეს მიდგომა უკეთ შეესაბამება კლინიკურ რეალობას, სადაც სისტემა ექიმის ჩანაცვლებას კი არ ცდილობს,
არამედ ეხმარება მას რისკის მიხედვით ინფორმაციის დალაგებაში.

\section{კვლევის ძლიერი მხარეები და შეზღუდვები}
ძლიერი მხარეებია მოდულური pipeline, leakage audit, deterministic seed protocol,
bootstrap სტატისტიკა და calibration/conformal safety შეფასება.
შეზღუდვებია ChestMNIST-ის 10\%-იანი extracted subset, frozen ELIXR-C ნიშნები,
გარე კლინიკური ვალიდაციის არქონა და დემოგრაფიული fairness audit-ის შეზღუდულობა.

\section{კლინიკური მნიშვნელობა}
მოდელი არ უნდა განიხილებოდეს ავტონომიურ დიაგნოსტიკურ პროდუქტად.
მისი სწორი გამოყენების სცენარია decision-support: საეჭვო პათოლოგიების prioritization,
დაბალკალიბრირებული ან გაურკვეველი შემთხვევების მონიშვნა და radiologist review-ის მხარდაჭერა.
კლინიკურ დანერგვამდე აუცილებელია prospective validation, ეთიკური თანხმობა,
PACS workflow-ში ტესტირება და post-deployment drift monitoring.

\iffalse
წინამდებარე თავი ეძღვნება მე-3 თავში წარმოდგენილი შედეგების ინტერპრეტაციას,
კრიტიკულ შეფასებას და მათი მნიშვნელობის განხილვას სამეცნიერო და კლინიკურ კონტექსტში.

\section{შედეგების ინტერპრეტაცია (Interpretation of Results)}

\subsection{რაოდენობრივი შედეგების მნიშვნელობა}
მე-3 თავში მიღებული შედეგები აჩვენებს, რომ EfficientNet-B0-ზე აგებული MultiLabelModel
მრავალნიშნიან ChestMNIST ამოცანაზე აღწევს პრაქტიკულად მნიშვნელოვან დისკრიმინაციულ ხარისხს.
მაკრო AUC-ROC = 0.7955 მიუთითებს, რომ მოდელი საშუალოდ კარგად განასხვავებს 14 პათოლოგიას,
თუმცა მაკრო F1 = 0.2413 ადასტურებს, რომ იშვიათ კლასებზე სწორი ზღვრის პოვნა კვლავ რთულია.

ამავე დროს, ბინარულ ჭრილში (\textit{ნებისმიერი პათოლოგია} vs \textit{არც ერთი})
AUC-ROC = 0.7418 და F1 = 0.6950 ცხადყოფს, რომ მოდელს უკეთ შეუძლია ზოგადი პათოლოგიურობის
იდენტიფიკაცია, ვიდრე თითოეული კონკრეტული ლეიბლის მაღალი სიზუსტით განცალკევება.
ეს განსხვავება ტიპურია დისბალანსირებული მრავალნიშნიანი ამოცანებისთვის,
სადაც უარყოფითი მაგალითების დომინანტობა ამცირებს მაკრო-F1-ს მიუხედავად კარგი AUC-ისა.

\subsection{ხარისხობრივი შედეგების და ვიზუალიზაციის ინტერპრეტაცია}
ხარისხობრივი ანალიზი (თავი 3.2) აჩვენებს, რომ მოდელი უფრო საიმედოდ მუშაობს იმ შემთხვევებში,
სადაც პათოლოგიური ნიშნები მკაფიოა, ხოლო შეცდომები უფრო ხშირად ჩნდება:

\begin{enumerate}[leftmargin=1.6em]
    \item დაბალკონტრასტულ ან სუსტი გამოვლინების შემთხვევებში;
    \item არტეფაქტებით დამახინჯებულ გამოსახულებებზე;
    \item ურთიერთმსგავსი რადიოლოგიური გამოვლინებების მქონე კლასებს შორის.
\end{enumerate}

Grad-CAM/attention ვიზუალიზაციები (თუ დამატებულია) განსაკუთრებით მნიშვნელოვანია კლინიკური ნდობისთვის:
სწორი პროგნოზებისას ყურადღების კონცენტრაცია ჩვეულებრივ ემთხვევა რელევანტურ ანატომიურ რეგიონებს,
ხოლო შეცდომით პროგნოზებში ხშირად ჩანს არარელევანტურ უბნებზე ფოკუსი.
ეს ხსნის იმას, თუ რატომ არის საჭირო მოდელის გადაწყვეტილების ექსპერტული გადამოწმება კლინიკურ გამოყენებამდე.

\subsection{კვლევის მიზნებისა და ამოცანების შესრულების შეფასება}
შესავალში დასმული ძირითადი ამოცანები შეფასდა შემდეგნაირად:

\begin{enumerate}[leftmargin=1.6em]
    \item \textbf{მონაცემთა ბაზის მომზადება}: შესრულებულია (ChestMNIST split + preprocessing).
    \item \textbf{მოდელის შერჩევა/იმპლემენტაცია}: შესრულებულია (MultiLabelModel, EfficientNet-B0 backbone).
    \item \textbf{სწავლება და ვალიდაცია}: შესრულებულია (ოპტიმიზებული ზღურბლები, early stopping პროტოკოლი).
    \item \textbf{საბოლოო შეფასება}: შესრულებულია (ხელშეუხებელი test split, ძირითადი მეტრიკებით).
    \item \textbf{შედეგების კრიტიკული ანალიზი}: შესრულებულია (რაოდენობრივი + ხარისხობრივი + სტატისტიკური შეზღუდვები).
\end{enumerate}

შესაბამისად, კვლევის ძირითადი მიზანი მიღწეულია:
შემუშავებულია და შეფასებულია ეფექტური საბაზისო მრავალნიშნიანი მოდელი,
რომელიც შეიძლება გამოყენებულ იქნას შემდგომი გაუმჯობესების პლატფორმად.

\section{შედარება არსებულ კვლევებთან (Comparison with Existing Studies)}

\subsection{თანხვედრა და განსხვავებები ეფექტურობის მაჩვენებლებში}
მიღებული შედეგები (AUC-ROC = 0.7955, მაკრო-F1 = 0.2413) თავსებადია იმ ფაქტთან,
რომ ChestMNIST წარმოადგენს რთულ, დისბალანსირებულ მრავალნიშნიან ბენჩმარკს
\cite{yang2023medmnistv2,wang2017chestxray14}.

სხვა პუბლიკაციებთან პირდაპირი რიცხვითი შედარება სიფრთხილით უნდა ჩატარდეს,
რადგან ხშირად განსხვავდება:

\begin{enumerate}[leftmargin=1.6em]
    \item მონაცემთა გაყოფა (patient-level vs image-level);
    \item preprocessing/augmentation მილსადენი;
    \item შეფასების პროტოკოლი (threshold tuning, averaging policy);
    \item გამოყენებული მეტრიკების კომპოზიცია.
\end{enumerate}

ამიტომ ამ ნაშრომის შედეგები უნდა განიმარტოს როგორც კონკრეტულ, მკაფიოდ აღწერილ პროტოკოლში
მიღებული რეპროდუცირებადი შედეგები.

\subsection{ჩვენი მიდგომის პოტენციური უპირატესობები}
MultiLabelModel-ის გამოყენებულ კონფიგურაციას აქვს რამდენიმე პრაქტიკული უპირატესობა:

\begin{enumerate}[leftmargin=1.6em]
    \item \textbf{ინჟინერული სიმარტივე}: EfficientNet-B0 კომპაქტურია და ადვილად ტრენირებადია.
    \item \textbf{ტრანსფერული სწავლების ეფექტი}: pretrained წონები ამცირებს კონვერგენციის რისკს მცირე/ხმაურიან მონაცემებზე.
    \item \textbf{კლასთა დისბალანსზე მორგება}: ოპტიმიზებული ზღურბლები და შესაბამისი loss დიზაინი ზრდის პრაქტიკულ გამოსადეგობას.
    \item \textbf{რეპროდუცირებადობა}: ფიქსირებული split და მკაფიო შეფასების მილსადენი ამარტივებს შედეგების გადამოწმებას.
\end{enumerate}

\subsection{ჩვენი მიდგომის შეზღუდვები არსებულ კვლევებთან შედარებით}
არსებული ტოპ-დონის კვლევები ხშირად იყენებს:
(ა) უფრო დიდ და მულტიცენტრულ მონაცემთა ნაკრებებს,
(ბ) მრავალმოდელურ ანსამბლებს,
(გ) გარე ვალიდაციას დამოუკიდებელ კოჰორტებზე.
წინამდებარე კვლევა ამ ეტაპზე ფოკუსირებულია ერთ ძირითად მოდელზე,
რაც აძლიერებს კონტროლირებადობას, მაგრამ ზღუდავს აბსოლუტური SOTA შედარების საშუალებას.

\section{კლინიკური მნიშვნელობა და პოტენციური გამოყენება (Clinical Significance and Potential Applications)}

\subsection{პოტენციური კლინიკური სარგებელი}
მიღებული შედეგები მიუთითებს რამდენიმე პრაქტიკულ გამოყენებაზე:

\begin{enumerate}[leftmargin=1.6em]
    \item \textbf{ტრიაჟის მხარდაჭერა}: პათოლოგიური შემთხვევების პრიორიტეტიზაცია სწრაფი რეაგირებისთვის;
    \item \textbf{მეორე მკითხველი}: რადიოლოგის გადაწყვეტილების მხარდაჭერა და გამოტოვების რისკის შემცირება;
    \item \textbf{სტანდარტიზაცია}: ინტერ-ოპერატორული ვარიაბელობის ნაწილობრივი შემცირება;
    \item \textbf{სასწავლო ღირებულება}: რთულ შემთხვევებზე ყურადღების ვიზუალიზაცია ახალგაზრდა სპეციალისტებისთვის.
\end{enumerate}

\subsection{პოტენციური გამოყენების სცენარები}
კლინიკურ პროცესში მოდელი შეიძლება ინტეგრირდეს როგორც:

\begin{enumerate}[leftmargin=1.6em]
    \item PACS/RIS გარემოში წინასწარი სკორინგის მოდული;
    \item სამუშაო სიის რისკზე დაფუძნებული პრიორიტეტიზატორი;
    \item ხარისხის კონტროლის ალგორითმი რეტროსპექტული აუდიტისთვის.
\end{enumerate}

\subsection{დანერგვამდე აუცილებელი ნაბიჯები}
კლინიკურ გამოყენებამდე საჭიროა:

\begin{enumerate}[leftmargin=1.6em]
    \item დამოუკიდებელი გარე ვალიდაცია (CheXpert/MIMIC-CXR);
    \item პროსპექტული კვლევა რეალურ workflow-ში;
    \item ინსტიტუციური ეთიკური და მარეგულირებელი შეფასება;
    \item უწყვეტი მონიტორინგი drift-ის, calibration-ისა და failure-mode-ებისთვის.
\end{enumerate}

\section{კვლევის ძლიერი მხარეები (Strengths of the Study)}
კვლევის ძირითადი ძლიერი მხარეებია:

\begin{enumerate}[leftmargin=1.6em]
    \item კლინიკურად აქტუალური ამოცანის არჩევა (14-კლასიანი მრავალნიშნიანი CXR კლასიფიკაცია);
    \item მკაფიო და გამჭვირვალე ექსპერიმენტული პროტოკოლი;
    \item ხელშეუხებელ სატესტო ნაკრებზე საბოლოო შეფასება;
    \item მრავალმხრივი მეტრიკების გამოყენება (AUC, F1, binary vs multilabel);
    \item ხარისხობრივი ანალიზის ჩართვა შეცდომების ტიპოლოგიის გასაგებად;
    \item პრაქტიკულად გამოყენებადი და რესურსულად რეალისტური არქიტექტურა;
    \item რეპროდუცირებადობაზე ორიენტირებული კვლევითი დიზაინი.
\end{enumerate}

\section{კვლევის შეზღუდვები (Limitations of the Study)}

\subsection{მონაცემთა ბაზის შეზღუდვები}
ChestMNIST-ის სტანდარტიზებული ფორმატი კვლევისთვის მოსახერხებელია,
თუმცა მისი გამოყენება ქმნის შემდეგ შეზღუდვებს:

\begin{enumerate}[leftmargin=1.6em]
    \item დისბალანსი კლასებს შორის;
    \item შეზღუდული დემოგრაფიული/კლინიკური მეტამონაცემები;
    \item გარე ვალიდაციის საჭიროება სხვადასხვა ინსტიტუციურ დომენზე.
\end{enumerate}

\subsection{მოდელის შეზღუდვები}
EfficientNet-B0 არის ძლიერი საბაზისო მოდელი, მაგრამ:

\begin{enumerate}[leftmargin=1.6em]
    \item იშვიათ კლასებზე მაკრო-F1 კვლავ დაბალია;
    \item არ არის გამოყენებული სრული ანსამბლური/მულტიმოდალური სტრატეგია;
    \item ინტერპრეტირებადობის მეთოდები (მაგ., Grad-CAM) სრულ მიზეზობრივ ახსნას არ იძლევა.
\end{enumerate}

\subsection{მეთოდოლოგიური შეზღუდვები}
კვლევა რეტროსპექტულია და არ მოიცავს პროსპექტულ კლინიკურ ინტეგრაციას.
ასევე, ჰიპერპარამეტრების სივრცე პრაქტიკულად ამოწურული არ არის,
რაც ნიშნავს, რომ უკეთესი კონფიგურაციის პოვნის პოტენციალი რჩება.

\subsection{შეჯამება}
შეზღუდვების მიუხედავად, მიღებული შედეგები ადასტურებს,
რომ მოცემულ კონფიგურაციაში MultiLabelModel წარმოადგენს მყარ საბაზისო სისტემას,
რომელიც მზადაა შემდეგ ეტაპზე გადასასვლელად: გარე ვალიდაცია, ინტერპრეტირებადობის გაძლიერება და
კლინიკურ workflow-ში ეტაპობრივი ინტეგრაცია.

\fi

\chapter{იმპლემენტაციის დეტალური ანალიზი}

\section{პროგრამული არქიტექტურის საერთო ხედვა}
CXR-Synapse-ის პროგრამული იმპლემენტაცია აგებულია კვლევითი pipeline-ის პრინციპით:
მონაცემთა მომზადება, feature extraction, მოდელის სწავლება, შეფასება, სტატისტიკური ანალიზი
და აუდიტი ერთმანეთისგან განცალკევებულია. ასეთი სტრუქტურა აუცილებელია სამაგისტრო კვლევისთვის,
რადგან თითოეული ეტაპის დამოუკიდებლად გადამოწმება შესაძლებელია.
კოდის ძირითადი ფაილებია \texttt{data.py}, \texttt{model.py}, \texttt{train.py},
\texttt{evaluators.py}, \texttt{utils.py}, \texttt{audit.py} და \texttt{main.py}.

\texttt{main.py} ასრულებს ორკომპონენტიან როლს. პირველი ნაწილი იწყებს სწავლების პროცესს:
იტვირთება embedding dataset-ები, იანგარიშება co-occurrence adjacency matrix,
მზადდება RadLex embedding-ები და იწყება სამწევრიანი ensemble-ის სწავლება.
მეორე ნაწილი განახორციელებს შეფასებას: ირჩევა ვალიდაციაზე class-specific threshold-ები,
იტვირთება single-model baseline, ეშვება deep ensemble MC-Dropout inference,
ითვლება bootstrap confidence interval, paired p-value და conformal coverage.
ამით მთავარი ფაილი ინარჩუნებს კვლევის ლოგიკურ სურათს და არ მალავს რომელიმე მნიშვნელოვან ნაბიჯს.

\section{მონაცემთა ფენა}
\texttt{data.py}-ში მონაცემები წარმოდგენილია ორი განსხვავებული სახით.
პირველი არის raw ChestMNIST dataset, რომელიც გამოიყენება leakage audit-ისთვის და hashing-ისთვის.
მეორე არის \texttt{EmbeddingDataset}, რომელიც კითხულობს უკვე ექსტრაქტირებულ ნიშნებს.
ასეთი გაყოფა პრაქტიკულია: ELIXR-C/CXR-Foundation მოდელით feature extraction ძვირია,
ხოლო downstream classifier-ის კვლევა ბევრად სწრაფად მეორდება precomputed embedding-ებით.

მონაცემთა integrity check მნიშვნელოვანი დაცვის მექანიზმია.
თუ embedding ფაილში features და labels რაოდენობა ერთმანეთს არ ემთხვევა,
კოდი ორივე მასივს საერთო მინიმალურ ზომამდე ჭრის და structured log-ში აფიქსირებს warning-ს.
სამეცნიერო თვალსაზრისით ეს ნაბიჯი მნიშვნელოვანია, რადგან feature-label mismatch უმნიშვნელო
ტექნიკურ შეცდომად შეიძლება გამოიყურებოდეს, მაგრამ სინამდვილეში მთლიან ექსპერიმენტს აბათილებს.

\section{Feature extraction ფენა}
\texttt{Extraction/extraction\_pipeline.py} მუშაობს resumable extraction პრინციპით.
ფაილი ჯერ ქმნის TFRecord cache-ს 10\%-იანი subset-ისთვის, შემდეგ კი TensorFlow SavedModel-ით
იღებს feature representation-ს. checkpoint-ები ინახება batch-ებად, რაც გრძელვადიანი პროცესის
შეწყვეტის შემთხვევაში მუშაობის გაგრძელების საშუალებას იძლევა.
ეს განსაკუთრებით მნიშვნელოვანია დიდი სამედიცინო მოდელების გამოყენებისას, სადაც ერთი extraction run
შეიძლება საათობით გაგრძელდეს.

ფინალური \texttt{.pt} ფაილი შეიცავს features, labels, hashes, model version-ს და extraction date-ს.
ეს ქმნის provenance chain-ს: შესაძლებელია დადგინდეს, რომელი backbone-ის რომელი extraction ვერსიით
არის შექმნილი downstream training data. სამაგისტრო ნაშრომისთვის ასეთი provenance მნიშვნელოვანია,
რადგან კვლევის შედეგი არ უნდა იყოს მხოლოდ ერთი ლოკალური გაშვების აღწერა;
ის უნდა იყოს ხელახლა აღდგენადი პროცესი.

\section{მოდელის ფენა}
\texttt{model.py}-ში CXR-Synapse წარმოდგენილია \texttt{CXR\_Synapse\_Foundation} კლასით.
მისი პირველი კომპონენტი არის dimensionality reduction block, რომელიც 1376-განზომილებიან
ELIXR-C feature-ს 384-განზომილებიან სივრცეში გადაიყვანს.
შემდეგ ემატება 2D sine-cosine positional embedding, რომელიც სივრცითი grid-ის ორგანიზაციას ინარჩუნებს.

მეორე მნიშვნელოვანი კომპონენტია \texttt{PathologyCrossAttention}.
აქ თითოეული პათოლოგია მოქმედებს როგორც query, ხოლო ვიზუალური patch-ები როგორც key/value.
კლინიკურად ეს ნიშნავს, რომ მოდელი თითოეული დაავადებისთვის ცალკე ეძებს შესაბამის ვიზუალურ მტკიცებულებას,
ნაცვლად იმისა, რომ მთლიანი გამოსახულების ერთი global vector პირდაპირ 14 sigmoid logit-ად გადააქციოს.
Graph message passing query-ებს აცნობს თანაარსებობის სტრუქტურას, რის გამოც pneumonia,
infiltration, effusion და სხვა დაკავშირებული ნიშნები სრულიად დამოუკიდებლად აღარ მუშავდება.

\section{სწავლების ფენა}
\texttt{train.py} აერთიანებს deterministic seed setting-ს, RadLex embedding loading-ს,
logit adjustment-ს, evidential loss-ს, LISA anchor loss-ს და SWA-ს.
სასწავლო ციკლი თითოეულ seed-ზე ცალ-ცალკე სრულდება, რის შედეგადაც იქმნება checkpoint-ები
\texttt{CXR\_Synapse\_Foundation\_Seed\_42.pth}, \texttt{Seed\_43.pth} და \texttt{Seed\_44.pth}.

LISA loss-ის როლია სემანტიკური topology-ის შენარჩუნება.
თუ training მხოლოდ label prediction loss-ს დაეყრდნობოდა, disease query-ები შეიძლება ისეთი მიმართულებით
გადაწეულიყო, რომელიც მოკლევადიანად აუმჯობესებს ვალიდაციის AUC-ს, მაგრამ შლის RadLex/BioViL-T-ის
კლინიკურ სემანტიკას. Anchor loss ამ drift-ს ზღუდავს და query space-ს უბრუნებს სამედიცინო ტერმინების
საწყის კავშირს.

\section{შეფასების ფენა}
\texttt{evaluators.py}-ში ორი შეფასების რეჟიმია.
\texttt{validate} გამოიყენება training loop-ში და აბრუნებს AUC, F1, mAP, raw probabilities და labels-ს.
\texttt{DeepEnsembleTTAEvaluator} კი საბოლოო შეფასებისთვის გამოიყენება.
მიუხედავად იმისა, რომ კლასის სახელში TTA წერია, მიმდინარე იმპლემენტაციის ძირითადი uncertainty წყაროა
ensemble წევრების საშუალო და MC-Dropout pass-ები; გამოსახულების horizontal flip augmentation
მიმდინარე feature-based pipeline-ში პირდაპირ არ გამოიყენება.

ამ განსხვავების ტექსტში დაფიქსირება მნიშვნელოვანია აკადემიური პატიოსნებისთვის.
ნაშრომი არ უნდა ამტკიცებდეს ისეთ კომპონენტს, რომელიც კოდში არ არსებობს.
სწორედ ამიტომ შედეგების თავში emphasis გადატანილია deep ensemble + MC-Dropout-ზე,
ხოლო Test-Time Augmentation არ არის წარმოდგენილი როგორც ფაქტობრივად შესრულებული ეტაპი.

\chapter{მონაცემთა აუდიტი და რეპროდუცირებადობა}

\section{Leakage-ის კონტროლი}
სამედიცინო AI კვლევებში leakage ერთ-ერთი ყველაზე სახიფათო პრობლემაა.
თუ ერთი და იგივე პაციენტის ან ერთი და იგივე გამოსახულების ვერსია აღმოჩნდება როგორც train,
ისე test ნაწილში, მოდელი რეალურად არ ზოგადდება, არამედ იმახსოვრებს.
CXR-Synapse pipeline იყენებს raw pixel hash-ებს, რათა იპოვოს bit-level duplicate-ები.
ვალიდაციიდან იშლება train-ში არსებული hash-ები, ხოლო test მოწმდება train და cleaned-val hash-ების
ერთიან reference set-თან.

\section{Forensic audit}
\texttt{audit.py} მონაცემთა ხარისხს რამდენიმე კუთხით ამოწმებს.
information-theoretic ნაწილი ითვლის mutual information-ს intensity statistics-სა და labels-ს შორის.
spatial morphology ნაწილი ქმნის global mean image-ს, pixel-wise variance-ს და edge-density map-ს.
clinical logic ნაწილი ითვლის label cardinality-ს, phi coefficient-ს, Cramer's V-ს და პირობით prevalence-ს.
ამგვარი audit არ ცვლის კლინიკურ ექსპერტიზას, მაგრამ აჩვენებს, სად შეიძლება dataset-ში არსებობდეს
არაბუნებრივი კორელაცია ან label artifact.

\section{დეტერმინიზმის პროტოკოლი}
\texttt{config.py} აფიქსირებს seed-ს Python random, NumPy, PyTorch CPU/GPU და CUDA backend-ებისთვის.
ჩართულია deterministic algorithms warn-only რეჟიმში და \texttt{CUBLAS\_WORKSPACE\_CONFIG}.
ასეთი მკაცრი რეჟიმი ზოგჯერ ამცირებს სიჩქარეს, მაგრამ ზრდის გამეორებადობას.
სამაგისტრო ნაშრომისთვის ეს tradeoff გამართლებულია, რადგან კვლევის შედეგი უნდა იყოს
არა მხოლოდ მაღალი ქულა, არამედ გამჭვირვალე და ხელახლა შესამოწმებელი პროცესი.

\section{სტატისტიკური დასაბუთება}
სატესტო AUC-ის confidence interval ითვლება bootstrap sampling-ით.
ყოველ გამეორებაზე test set ირჩევა replacement-ით, შემდეგ ითვლება metric.
2.5 და 97.5 percentile ქმნის 95\%-იან ნდობით ინტერვალს.
paired bootstrap test კი ერთსა და იმავე resample-ზე ადარებს ensemble prediction-ს single-model baseline-ს.
ეს სწორია, რადგან ორი მოდელი ერთსა და იმავე მაგალითებზე ფასდება და მათი შეცდომები დამოუკიდებელი არ არის.

\section{შედეგების provenance}
საბოლოო შედეგების სრული გამეორებისთვის საჭიროა:
\begin{enumerate}[leftmargin=1.6em]
    \item extraction cache და embedding ფაილები;
    \item სამი seed checkpoint;
    \item RadLex embedding cache;
    \item validation probabilities threshold optimization-ისთვის;
    \item test probabilities bootstrap და conformal analysis-ისთვის;
    \item code snapshot და requirements.
\end{enumerate}
ამ ელემენტების გარეშე PDF-ში მოცემული შედეგები მხოლოდ ტექსტურ ანგარიშად დარჩება.
ამ ელემენტებით კი ნაშრომი გარდაიქმნება სრულ კვლევით არტეფაქტად.

\chapter{კლინიკური, ეთიკური და სამართლებრივი განხილვა}

\section{Decision-support, არა ავტონომიური დიაგნოზი}
CXR-Synapse-ის გამოყენება სწორია მხოლოდ decision-support ჩარჩოში.
მოდელის output არის probability score და prediction set, არა საბოლოო დიაგნოზი.
რადიოლოგმა უნდა დაინახოს, რომ პროგნოზი არის AI-generated და საჭიროების შემთხვევაში უგულებელყოს იგი.
ეს განასხვავებს კვლევით მოდელს სამედიცინო მოწყობილობისგან.

\section{კალიბრაციის მნიშვნელობა კლინიკაში}
კლინიკურ workflow-ში ცუდად კალიბრირებული მაღალი probability განსაკუთრებით სახიფათოა.
თუ მოდელი ხშირად აძლევს 0.95 score-ს ისეთ შემთხვევებს, რომლებიც რეალურად გაცილებით ნაკლებად
დადასტურებულია, ექიმი შეიძლება ზედმეტად ენდოს სისტემას.
ECE=0.0120 ამ რისკს ამცირებს, მაგრამ საბოლოო დასკვნისთვის საჭიროა გარე ვალიდაცია.
შიდა ChestMNIST calibration არ ნიშნავს, რომ სხვა საავადმყოფოს მონაცემებზე იგივე calibration შენარჩუნდება.

\section{Conformal safety-ის როლი}
Conformal prediction-ის მთავარი უპირატესობა ის არის, რომ იგი მოდელს აძლევს ფრთხილი პასუხის ფორმას.
თუ შემთხვევა მარტივია, prediction set შეიძლება მცირე იყოს.
თუ შემთხვევა რთულია, set ფართოვდება და ექიმს აფრთხილებს, რომ რამდენიმე პათოლოგია უნდა განიხილოს.
ეს უკეთესია, ვიდრე იძულებითი ერთი label, განსაკუთრებით მაშინ, როცა CXR გამოსახულებები მრავალნიშნიანია
და ერთი პაციენტი რამდენიმე პათოლოგიას შეიძლება ატარებდეს.

\section{გარე ვალიდაციის აუცილებლობა}
ChestMNIST არის სასარგებლო benchmark, მაგრამ იგი არ არის საკმარისი კლინიკური დანერგვისთვის.
საჭიროა CheXpert, MIMIC-CXR ან ადგილობრივი hospital cohort, სადაც label policy,
scanner distribution, patient demographics და acquisition protocol განსხვავებულია.
გარე ვალიდაციამ უნდა გაზომოს არა მხოლოდ AUC, არამედ sensitivity/specificity, calibration,
coverage, subgroup performance და workflow impact.

\section{ეთიკური შეზღუდვები}
ნებისმიერი სამედიცინო AI სისტემა შეიძლება შეიცავდეს bias-ს.
თუ training data-ში გარკვეული ჯგუფები არასაკმარისადაა წარმოდგენილი, მოდელმა შეიძლება მათზე უარესი შედეგი აჩვენოს.
ამ ნაშრომის ფარგლებში სრული demographic audit შეუძლებელია, რადგან ChestMNIST შესაბამის metadata-ს არ იძლევა.
ამიტომ სამომავლო სამუშაოში fairness analysis უნდა გახდეს ერთ-ერთი მთავარი მოთხოვნა.

\section{პერსონალური მონაცემები და უსაფრთხოება}
მიუხედავად იმისა, რომ MedMNIST საჯარო და დეიდენტიფიცირებული benchmark-ია,
კლინიკურ გარემოში ნებისმიერი ახალი მონაცემი უნდა დამუშავდეს პერსონალური მონაცემების დაცვის წესებით.
საჭიროა access control, audit logging, encrypted storage და მკაფიო data retention policy.
მოდელის retraining pipeline არ უნდა ინახავდეს პაციენტის იდენტიფიცირებად ინფორმაციას.

\section{დანერგვის წინაპირობები}
რეალურ კლინიკაში გადასვლამდე აუცილებელია:
\begin{enumerate}[leftmargin=1.6em]
    \item prospective blinded validation;
    \item radiologist reader study;
    \item failure mode and effects analysis;
    \item PACS integration-ის უსაფრთხო ტესტირება;
    \item post-deployment monitoring dashboard;
    \item პროცედურა, რომელიც განსაზღვრავს ვინ აგებს პასუხს AI alert-ის იგნორირების ან შეცდომის შემთხვევაში.
\end{enumerate}

\chapter{აბლაციური კვლევის და დამატებითი ექსპერიმენტების გეგმა}

\section{აბლაციური კვლევის მიზანი}
მიმდინარე შედეგები აჩვენებს CXR-Synapse-ის ერთიან ეფექტურობას, მაგრამ სამეცნიერო დასაბუთებისთვის
საჭიროა გაირკვეს, რომელი კომპონენტი რა წვლილს ქმნის. აბლაციური კვლევა საშუალებას იძლევა
მოდელის არქიტექტურა არ შეფასდეს როგორც შავი ყუთი. თითოეული მოდულის გამორთვით ან ჩანაცვლებით
იზომება, რამდენად იცვლება AUC, F1, ECE, conformal coverage და average set size.

აბლაცია უნდა ჩატარდეს იდენტურ მონაცემებზე, იდენტური seed-ებით და იგივე evaluation protocol-ით.
სხვა შემთხვევაში განსხვავება შეიძლება მოდულის ეფექტი კი არა, შემთხვევითი initialization-ის ან
სხვა training noise-ის შედეგი იყოს. რეკომენდებულია ყველა კონფიგურაციისთვის მინიმუმ სამი seed,
რათა საშუალო და standard deviation შეფასდეს.

\section{საწყისი baseline}
პირველი baseline უნდა იყოს მარტივი multilabel classifier, რომელიც იღებს ELIXR-C feature-ების
global pooled representation-ს და იყენებს independent sigmoid head-ს. ეს baseline პასუხობს კითხვას:
რამდენს გვაძლევს foundation representation მარტო, label graph-ისა და PaQ attention-ის გარეშე.
თუ CXR-Synapse მხოლოდ ამ baseline-ს ოდნავ აჯობებს, მაშინ არქიტექტურული სირთულე ნაკლებად გამართლებულია.
თუ სხვაობა ჩანს კალიბრაციასა და conformal coverage-ში, მაშინ წვლილი მხოლოდ AUC-ით არ უნდა შეფასდეს.

\section{PaQ attention-ის აბლაცია}
Pathology-as-Query attention-ის გამორთვისას disease-specific query-ები უნდა ჩანაცვლდეს shared classifier head-ით.
ასეთი აბლაცია აჩვენებს, რამდენად მნიშვნელოვანია თითოეული პათოლოგიისთვის ცალკე ვიზუალური მტკიცებულების ძებნა.
მოსალოდნელია, რომ PaQ განსაკუთრებით დაეხმარება იმ კლასებს, რომელთა ვიზუალური ნიშნები ლოკალიზებულია
ან რომელთა global image-level representation ადვილად იკარგება.

შეფასებისას მნიშვნელოვანია per-class breakdown.
მაგალითად, cardiomegaly უფრო global morphology-ს ეფუძნება, ხოლო nodule ან mass შეიძლება უფრო ლოკალური იყოს.
თუ PaQ ლოკალურ პათოლოგიებზე მეტ გაუმჯობესებას აჩვენებს, ეს გააძლიერებს არქიტექტურული არჩევანის
კლინიკურ ინტერპრეტაციას.

\section{Graph adjacency-ის აბლაცია}
Graph-guided query fusion-ის გამორთვა შესაძლებელია adjacency matrix-ის identity matrix-ით ჩანაცვლებით.
ამ შემთხვევაში query-ები მხოლოდ საკუთარ RadLex embedding-ს ეყრდნობა და აღარ იღებს co-occurrence შეტყობინებას.
შედარება აჩვენებს, რამდენად მნიშვნელოვანია pathology dependencies.

გრაფის ეფექტი განსაკუთრებით უნდა გამოჩნდეს თანაარსებობის მქონე label-ებზე.
თუ pneumonia და infiltration, effusion და atelectasis ან edema და cardiomegaly ხშირად ერთად გვხვდება,
graph message passing შესაძლოა მათი ერთობლივი pattern-ის დაჭერას დაეხმაროს.
თუმცა არსებობს რისკიც: თუ გრაფი ზედმეტად ძლიერია, იგი შეიძლება ასწავლიდეს მოდელს dataset-specific correlation-ს,
რომელიც სხვა კლინიკურ გარემოში აღარ იმუშავებს. ამიტომ graph ablation არა მხოლოდ performance-ის,
არამედ generalization-ის საკითხიცაა.

\section{LISA regularization-ის აბლაცია}
LISA anchor loss-ის გამორთვა ამოწმებს, რამდენად იცავს სემანტიკური topology training drift-ისგან.
თუ LISA-ს გარეშე AUC იზრდება, მაგრამ ECE უარესდება ან per-class stability მცირდება,
ეს მიუთითებს, რომ სუფთა discrimination objective შეიძლება ზედმეტად მოერგოს validation distribution-ს.
თუ LISA-ს დამატება ამცირებს drift-ს და აუმჯობესებს calibration-ს, იგი გამართლებულია როგორც safety-oriented
regularizer.

LISA-ს გავლენის გასაზომად რეკომენდებულია query topology distance-ის მონიტორინგი epoch-ების მიხედვით.
საწყისი RadLex cosine matrix და მიმდინარე query cosine matrix შეიძლება შედარდეს Frobenius norm-ით.
ეს მოგვცემს არა მხოლოდ საბოლოო metric-ს, არამედ დინამიკურ სურათს, როგორ იცვლება სემანტიკური სივრცე სწავლებისას.

\section{Evidential loss-ის აბლაცია}
Strictly Proper Beta Evidential Loss უნდა შედარდეს ჩვეულებრივ BCEWithLogitsLoss-თან.
ამ აბლაციის მთავარი კითხვა არის calibration და uncertainty.
შესაძლებელია BCE-მ მსგავსი AUC აჩვენოს, მაგრამ უფრო მაღალი ECE ჰქონდეს.
სამედიცინო კონტექსტში ასეთი შედეგი ნიშნავს, რომ evidential loss ღირებულია,
თუნდაც ranking metric-ში დიდი სხვაობა არ იყოს.

ასევე საჭიროა annealing coefficient-ის კვლევა.
ძალიან სწრაფმა KL regularization-მა შეიძლება შეაფერხოს discriminative learning,
ხოლო ძალიან სუსტმა regularization-მა uncertainty estimation ეფექტი დაკარგოს.
ამიტომ სასურველია მინიმუმ სამი ვარიანტი: annealing 10 epoch, 20 epoch და 35 epoch.

\section{SWA და ensemble აბლაცია}
SWA უნდა შედარდეს ბოლო checkpoint-თან და best validation checkpoint-თან.
SWA ხშირად აუმჯობესებს generalization-ს, რადგან წონების საშუალო ფართო optimum-ისკენ გადადის.
თუმცა ყველა dataset-ზე ეს ავტომატურად არ ხდება, ამიტომ საჭიროა ცალკე შეფასება.

Ensemble აბლაციაში უნდა შედარდეს:
\begin{enumerate}[leftmargin=1.6em]
    \item ერთი მოდელი dropout-ის გარეშე;
    \item ერთი მოდელი MC-Dropout-ით;
    \item სამი seed ensemble;
    \item სამი seed ensemble MC-Dropout-ით.
\end{enumerate}
ასეთი მატრიცა განასხვავებს epistemic uncertainty-ის ორ წყაროს: seed diversity-ს და stochastic inference-ს.

\section{Conformal wrapper-ის sensitivity}
Conformal predictor-ისთვის საჭიროა $\alpha$-ს sensitivity analysis: 0.05, 0.10 და 0.15.
დაბალი $\alpha$ უფრო მაღალი coverage-ს მოითხოვს და, როგორც წესი, ზრდის prediction set size-ს.
კლინიკური თვალსაზრისით ეს არის უსაფრთხოება-ეფექტიანობის tradeoff.
ძალიან დიდი set ექიმს ნაკლებად ეხმარება, ხოლო ძალიან მცირე set შეიძლება გამოტოვებდეს მნიშვნელოვან პათოლოგიას.

\section{Threshold optimization-ის შედარება}
მიმდინარე კოდი threshold-ებს validation probability percentile grid-ზე არჩევს.
შედარებისთვის შესაძლებელია fixed 0.5 threshold, Youden-style criterion და F1-optimized threshold.
Macro F1 განსაკუთრებით მგრძნობიარეა იშვიათ კლასებზე, ამიტომ thresholding strategy პირდაპირ გავლენას ახდენს
საბოლოო დასკვნაზე. ნაშრომში უნდა დაფიქსირდეს, რომ threshold-ები არ ირჩევა test set-ზე,
რათა თავიდან იქნას აცილებული evaluation leakage.

\chapter{Threats to Validity და რისკების მართვა}

\section{Internal validity}
Internal validity ეხება იმას, არის თუ არა მიღებული შედეგი მართლაც მოდელის არქიტექტურის შედეგი.
პოტენციური საფრთხეებია random seed sensitivity, preprocessing mismatch, feature-label misalignment,
threshold tuning leakage და checkpoint selection bias.
CXR-Synapse ამ საფრთხეებს ნაწილობრივ პასუხობს deterministic seed protocol-ით, embedding integrity check-ით,
validation-only threshold optimization-ით და test set-ის მხოლოდ საბოლოო შეფასებისთვის გამოყენებით.

მიუხედავად ამისა, საჭიროა დამატებითი კონტროლი.
ყველა ძირითადი შედეგი უნდა ინახებოდეს raw prediction format-ში.
ეს საშუალებას იძლევა შემდგომში გადაითვალოს AUC, ECE, conformal coverage და bootstrap interval
კოდის ხელახლა გაშვების გარეშე. კვლევის დაცვისას ასეთი ფაილები შეიძლება წარმოდგენილი იყოს როგორც
reproducibility bundle-ის ნაწილი.

\section{External validity}
External validity უკავშირდება სხვა dataset-ზე ან სხვა კლინიკურ გარემოში შედეგების გადატანას.
ChestMNIST არის compact benchmark და არ ასახავს სრულ clinical heterogeneity-ს.
განსხვავებული scanner, hospital protocol, patient population და label policy შეიძლება მნიშვნელოვნად შეცვალოს
მოდელის performance.

ამ საფრთხის მართვის ერთადერთი სანდო გზა არის external validation.
CheXpert და MIMIC-CXR განსხვავებული label extraction პოლიტიკით და გამოსახულების ხარისხით გამოირჩევა.
თუ CXR-Synapse ამ dataset-ებზეც ინარჩუნებს calibration-ს და coverage-ს, მაშინ არგუმენტი
კლინიკური რელევანტურობის შესახებ უფრო ძლიერი იქნება.

\section{Construct validity}
Construct validity სვამს კითხვას: ნამდვილად ზომავს თუ არა არჩეული metric იმას, რაც გვაინტერესებს.
AUC ზომავს ranking quality-ს, მაგრამ არ ამბობს, რამდენად სწორია probability.
F1 დამოკიდებულია threshold-ზე და დისბალანსირებულ კლასებში ადვილად მერყეობს.
ECE ზომავს calibration-ს, მაგრამ binning strategy-ზეა დამოკიდებული.
Conformal coverage ზომავს true label-ის მოხვედრას prediction set-ში, მაგრამ დიდ set-ს შეიძლება
კლინიკური utility დაბალი ჰქონდეს.

ამიტომ ნაშრომი metrics-ს ერთობლივად განიხილავს.
მაღალი AUC და დაბალი ECE ერთად უფრო ძლიერი არგუმენტია, ვიდრე თითოეული ცალკე.
Coverage და set size ასევე ერთად უნდა იყოს წარმოდგენილი: coverage-ის ზრდა სასარგებლოა მხოლოდ მაშინ,
თუ set size უკონტროლოდ არ იზრდება.

\section{Conclusion validity}
Conclusion validity ეხება სტატისტიკური დასკვნის სანდოობას.
ერთი run-ის შედეგი არასაკმარისია, რადგან neural network training stochastic პროცესია.
სწორედ ამიტომ გამოიყენება სამი seed ensemble და bootstrap confidence interval.
paired bootstrap test უკეთესია დამოუკიდებელ t-test-ზე, რადგან ორი მოდელის prediction-ები ერთსა და იმავე
test examples-ზეა მიღებული.

საფრთხე მაინც რჩება: bootstrap assumes test set approximates target distribution.
თუ test set პატარაა ან დისტრიბუციულად შეზღუდულია, confidence interval შეიძლება ზედმეტად ოპტიმისტური იყოს.
ამიტომ სტატისტიკური დასკვნა უნდა იყოს მოკრძალებული და არ უნდა გადაიზარდოს კლინიკური დანერგვის
დაუყოვნებელ რეკომენდაციაში.

\section{Engineering risks}
Engineering risk მოიცავს dependency drift-ს, path mismatch-ს, GPU/CPU numerical differences-ს და framework version changes-ს.
\texttt{requirements.txt} ამცირებს ამ პრობლემას, მაგრამ სრული reproducibility-სთვის სასურველია lock file
ან container image.
Feature extraction pipeline იყენებს TensorFlow SavedModel-ს, ხოლო downstream training PyTorch-ზეა აგებული.
ორი framework-ის კომბინაცია მოქნილია, მაგრამ ზრდის environment complexity-ს.

\section{Clinical risks}
კლინიკური რისკები მოიცავს missed critical findings-ს, false alarm fatigue-ს, automation bias-ს და workflow delay-ს.
Conformal prediction და calibration ამ რისკებს ამცირებს, მაგრამ არ აქრობს.
საჭიროა user interface, რომელიც აჩვენებს uncertainty-ს გასაგები ფორმით და არ მალავს მოდელის შეზღუდვებს.
ასევე აუცილებელია audit trail: ვინ ნახა AI output, როდის, რა გადაწყვეტილება მიიღო და იყო თუ არა disagreement.

\chapter{კლინიკური ვალიდაციის სრული პროტოკოლი}

\section{კვლევის დიზაინი}
კლინიკური ვალიდაციის რეკომენდებული დიზაინია prospective, blinded reader study.
მოდელი წინასწარ უნდა გაიყინოს და validation protocol უნდა დარეგისტრირდეს ექსპერიმენტის დაწყებამდე.
რადიოლოგებმა არ უნდა იცოდნენ მოდელის prediction-ები პირველ შეფასების რაუნდში,
რათა baseline human performance ობიექტურად გაიზომოს.
მეორე რაუნდში შეიძლება მიეცეთ AI output და შეფასდეს workflow impact.

\section{კოჰორტის შერჩევა}
კოჰორტა უნდა მოიცავდეს მინიმუმ რამდენიმე ასეულ CXR კვლევას, სასურველია 500--1000 შემთხვევა.
ნიმუშში უნდა იყოს ნორმალური შემთხვევები, ხშირი პათოლოგიები და იშვიათი პათოლოგიები.
თუ rare class ძალიან ცოტაა, sensitivity-ის შეფასება არასტაბილური იქნება.
საჭიროა inclusion/exclusion criteria: ასაკი, გამოსახულების ხარისხი, projection type,
portable versus standard acquisition და პაციენტის repeated studies-ის მართვა.

\section{Reference standard}
Gold standard უნდა შეიქმნას მინიმუმ ორი დამოუკიდებელი რადიოლოგის შეფასებით.
disagreement შემთხვევაში საჭიროა adjudication მესამე senior radiologist-ის მიერ.
თუ label-ები მხოლოდ report mining-ით მიიღება, uncertainty მაღალია და validation conclusion სუსტდება.
Reference standard-ის ხარისხი პირდაპირ ზღუდავს მოდელის შეფასების ხარისხს.

\section{Primary და secondary endpoints}
Primary endpoints უნდა იყოს sensitivity და specificity clinically important label-ებისთვის,
მაგალითად pneumothorax, edema, consolidation და cardiomegaly.
Secondary endpoints შეიძლება იყოს macro AUC, per-class AUC, ECE, Brier score, conformal coverage,
average set size, reading time და radiologist confidence.

ასევე სასურველია workflow endpoint:
AI output ამცირებს თუ არა reading time-ს?
ზრდის თუ არა confidence-ს რთულ შემთხვევებში?
იწვევს თუ არა false alert fatigue-ს?
ეს კითხვები აუცილებელია, რადგან კარგი offline metric ყოველთვის არ ნიშნავს კარგ clinical utility-ს.

\section{Safety monitoring}
კვლევის დროს უნდა არსებობდეს safety monitoring plan.
ყველა missed critical finding უნდა ჩაიწეროს და გაანალიზდეს.
თუ მოდელი ხშირად გამოტოვებს სასწრაფო პათოლოგიას, deployment უნდა შეჩერდეს.
False positives ასევე მნიშვნელოვანია, რადგან ზედმეტი alerts კლინიცისტს შეიძლება მიაჩვიოს სისტემის იგნორირებას.

\section{Subgroup analysis}
თუ metadata ხელმისაწვდომია, performance უნდა დაიყოს ასაკის, სქესის, acquisition device-ის,
hospital site-ის და projection type-ის მიხედვით.
Subgroup disparity შეიძლება დაიმალოს aggregate AUC-ში.
მაგალითად, საერთო AUC მაღალი იყოს, მაგრამ portable ICU images-ზე performance მკვეთრად დაბალი.
ასეთი განსხვავება კლინიკური უსაფრთხოების საკითხია და არა მხოლოდ სტატისტიკური დეტალი.

\section{Deployment monitoring}
პილოტური დანერგვის შემდეგ საჭიროა dashboard, რომელიც აკონტროლებს:
\begin{enumerate}[leftmargin=1.6em]
    \item prediction distribution drift-ს;
    \item calibration drift-ს;
    \item conformal set size-ის ზრდას;
    \item per-class sensitivity trend-ს;
    \item system downtime-ს;
    \item user override rate-ს.
\end{enumerate}
თუ ECE ან set size იზრდება, ეს შეიძლება მიუთითებდეს distribution shift-ზე.
ასეთ შემთხვევაში საჭიროა revalidation და არა ავტომატური retraining.

\chapter{ოპერაციული პროცედურები და დაცვისთვის მზადყოფნა}

\section{კვლევითი გარემოს აღდგენის პროცედურა}
ნაშრომის პრაქტიკული ღირებულება მნიშვნელოვნად იზრდება, თუ კომისიის ან სხვა მკვლევრისთვის შესაძლებელია
კვლევითი გარემოს აღდგენა. პროცედურა იწყება Python გარემოს შექმნით და \texttt{requirements.txt}-ში
მითითებული პაკეტების დაყენებით. შემდეგ უნდა შემოწმდეს CUDA/PyTorch თავსებადობა,
რადგან training და inference შედეგები შეიძლება მცირედ განსხვავდებოდეს CPU და GPU რეჟიმებს შორის.

შემდეგი ნაბიჯია მონაცემთა არტეფაქტების განთავსება მოსალოდნელ ბილიკებზე.
\texttt{data.py}-ში \texttt{EMBEDDING\_DIR} მიუთითებს \texttt{/workspace/Extraction/cxr\_embeddings\_10percent}
საქაღალდეზე. თუ გარემო ლოკალურ Windows მანქანაზე ეშვება, ეს ბილიკი უნდა შეიცვალოს შესაბამის absolute path-ად.
ეს ცვლილება კვლევის შედეგს არ ცვლის, მაგრამ აუცილებელია კოდის გაშვებისთვის.

\section{Feature extraction-ის operational checklist}
Feature extraction-ის დაწყებამდე უნდა შემოწმდეს:
\begin{enumerate}[leftmargin=1.6em]
    \item ELIXR-C/CXR-Foundation SavedModel-ის მდებარეობა;
    \item TensorFlow და tensorflow-text ვერსიები;
    \item disk space checkpoint-ებისა და final embedding ფაილებისთვის;
    \item TFRecord cache-ის schema version;
    \item extraction log file, სადაც process interruption და resume მდგომარეობა ჩანს.
\end{enumerate}

თუ extraction გაწყდა, pipeline თავიდან არ იწყებს ყველა sample-ს.
იგი კითხულობს checkpoint ფაილებს, აგროვებს უკვე დამუშავებულ hash-ებს და აგრძელებს დარჩენილი ჩანაწერებიდან.
ეს თვისება მნიშვნელოვანია, რადგან medical foundation model inference ხანგრძლივია და interruption პრაქტიკაში
ხშირად ხდება, განსაკუთრებით shared GPU გარემოში.

\section{Training-ის operational checklist}
Training-ის დაწყებამდე საჭიროა embedding ფაილების structural check.
\texttt{EmbeddingDataset} უკვე ამოწმებს features და labels რაოდენობას, მაგრამ სასურველია დამატებით
დაიბეჭდოს feature shape და label prevalence.
თუ რომელიმე class მთლიანად ნულოვანია extracted subset-ში, macro AUC და F1 interpretation სუსტდება.

Training-ისას უნდა შეინახოს:
\begin{enumerate}[leftmargin=1.6em]
    \item თითოეული seed-ის checkpoint;
    \item validation AUC trajectory;
    \item loss trajectory;
    \item adjacency threshold;
    \item logit adjustment vector;
    \item RadLex embedding file hash.
\end{enumerate}
ეს ინფორმაცია საჭიროა იმის დასადგენად, training სტაბილურად მიმდინარეობდა თუ შედეგი ერთი შემთხვევითი
checkpoint-ის ეფექტია.

\section{Evaluation-ის operational checklist}
Evaluation მხოლოდ training-ის დასრულების შემდეგ უნდა შესრულდეს.
პირველ რიგში, validation set-ზე ირჩევა thresholds.
შემდეგ single model baseline ფასდება test set-ზე.
მხოლოდ ამის შემდეგ ეშვება ensemble evaluator.
ასეთი თანმიმდევრობა იცავს test set-ს tuning-ისგან.

საბოლოო \texttt{results\_df} უნდა შეინახოს CSV ან JSON ფორმატში.
ამჟამად \texttt{main.py} შედეგს ბეჭდავს console-ში, მაგრამ სამაგისტრო submission bundle-ისთვის უკეთესია
ფაილად შენახვაც. რეკომენდებულია \texttt{results/summary.csv}, \texttt{results/test\_predictions.npy}
და \texttt{results/test\_labels.npy}.

\section{დაცვის პრეზენტაციის სტრუქტურა}
დაცვის 15--20 წუთიანი პრეზენტაცია უნდა იყოს მკაფიო და არა გადატვირთული.
რეკომენდებული სტრუქტურაა:
\begin{enumerate}[leftmargin=1.6em]
    \item პრობლემა: რატომ არის CXR multilabel classification რთული;
    \item კვლევის მიზანი და novelty;
    \item მონაცემები და leakage audit;
    \item მოდელის არქიტექტურა: ELIXR-C features, RadLex queries, graph-guided PaQ, evidential loss;
    \item training და evaluation protocol;
    \item ძირითადი შედეგები: AUC, ECE, coverage;
    \item შეზღუდვები და გარე ვალიდაციის გეგმა;
    \item დასკვნა.
\end{enumerate}

პრეზენტაციაში განსაკუთრებული ყურადღება უნდა მიექცეს იმას, რომ მოდელი არ არის კლინიკური პროდუქტი.
ეს წინასწარ პასუხობს კომისიის ეთიკურ კითხვებს და აჩვენებს, რომ ავტორი აცნობიერებს სამედიცინო AI-ის
საფრთხეებს.

\section{მოსალოდნელი კითხვები დაცვაზე}
კომისიამ შეიძლება იკითხოს, რატომ არ მოხდა end-to-end training.
პასუხი არის, რომ კვლევის მიზანი არ არის backbone-ის ნულიდან ან სრულად fine-tuning რეჟიმში გადამზადება.
ნაშრომი იკვლევს, როგორ შეიძლება ძლიერი CXR foundation representation-ის გამოყენება downstream,
სემანტიკურად მდიდარი და uncertainty-aware classifier-ის შესაქმნელად.

შესაძლებელია კითხვა: რატომ არის AUC 0.7913 საკმარისი?
აქ პასუხი უნდა იყოს ფრთხილი. AUC არ არის კლინიკური deployment-ის საკმარისი საფუძველი.
ნაშრომის წვლილი არის მთლიან ჩარჩოში: calibrated prediction, conformal coverage, leakage audit და
reproducible pipeline. მაღალი AUC-ის ნაცვლად აქ მთავარი არგუმენტი trustworthy evaluation-ია.

კიდევ ერთი კითხვა შეიძლება ეხებოდეს 10\%-იან subset-ს.
პასუხი: subset გამოყენებულია გამოთვლითი შეზღუდვის გამო, მაგრამ pipeline აგებულია ისე,
რომ სრული dataset-ის გამოყენება მხოლოდ extraction/training დროის ზრდას მოითხოვს და არა მეთოდოლოგიის შეცვლას.
ეს უნდა ჩაითვალოს შეზღუდვად და სამომავლო სამუშაოდ.

\section{კომისიის შეფასების კრიტერიუმებთან შესაბამისობა}
საკვლევი პრობლემის აქტუალობა ჩანს სამედიცინო imaging AI-ის მოთხოვნებიდან:
მრავალნიშნიანი CXR classification, calibration და uncertainty რეალურ კლინიკურ workflow-ში მნიშვნელოვანია.
პრაქტიკული მნიშვნელობა მდგომარეობს decision-support pipeline-ში, რომელიც ექიმის ჩანაცვლებას არ ცდილობს.
თეორიული მნიშვნელობა არის semantic label queries, co-occurrence graph და evidential/conformal logic-ის
ერთიან მოდელში გაერთიანება.

მეთოდოლოგიის შესაბამისობა დასმულ პრობლემასთან უზრუნველყოფილია იმით, რომ ყველა მეთოდი პირდაპირ უკავშირდება
ერთ კონკრეტულ გამოწვევას:
\begin{enumerate}[leftmargin=1.6em]
    \item foundation features პასუხობს მცირე dataset-ისა და representation learning-ის პრობლემას;
    \item graph-guided PaQ პასუხობს label dependency-ს;
    \item logit adjustment პასუხობს class imbalance-ს;
    \item evidential loss და ECE პასუხობს calibration-ს;
    \item conformal prediction პასუხობს safety coverage-ს;
    \item bootstrap პასუხობს სტატისტიკური სანდოობის მოთხოვნას.
\end{enumerate}

\section{ნაშრომის შესაძლო გაუმჯობესებები submission-მდე}
საბოლოო submission-მდე სასურველია დაემატოს რეალური ფიგურები.
პირველი ფიგურა უნდა იყოს pipeline diagram: raw ChestMNIST, ELIXR-C extraction, embedding dataset,
CXR-Synapse model, ensemble evaluation და conformal wrapper.
მეორე ფიგურა უნდა იყოს reliability diagram ან per-class ECE plot.
მესამე ფიგურა შეიძლება იყოს co-occurrence adjacency heatmap.
მეოთხე ფიგურა სასურველია აჩვენებდეს conformal set size distribution-ს.

ასევე სასურველია ცხრილში დაემატოს per-class prevalence.
ეს კომისიას დაეხმარება გაიგოს, რატომ არის macro F1 რთული metric და რატომ არის rare pathology recall
მნიშვნელოვანი, თუნდაც aggregate AUC მისაღები იყოს.

\section{საბოლოო არტეფაქტების სია}
სამაგისტრო ნაშრომთან ერთად სასურველია ინახებოდეს:
\begin{enumerate}[leftmargin=1.6em]
    \item PDF ფაილი;
    \item \LaTeX{} source;
    \item code snapshot;
    \item requirements;
    \item final checkpoints;
    \item prediction arrays;
    \item result summary;
    \item figure generation scripts;
    \item plagiarism/Turnitin report, უნივერსიტეტის პროცედურის შესაბამისად.
\end{enumerate}

\section{დაცვისას არგუმენტაციის ლოგიკა}
დაცვისას ავტორმა უნდა აჩვენოს არა მხოლოდ ის, რომ მოდელი მუშაობს, არამედ ისიც,
რომ კვლევითი არჩევანი თანმიმდევრულია. არგუმენტაცია შეიძლება აშენდეს შემდეგი ხაზით:
ChestMNIST მრავალნიშნიანი CXR ამოცანაა; ამ ამოცანაში პრობლემაა representation, imbalance,
label dependency და uncertainty; CXR-Synapse თითოეულ პრობლემას კონკრეტული მოდულით პასუხობს.
ეს ლოგიკა ეხმარება კომისიას დაინახოს, რომ არქიტექტურა შემთხვევითი კომპონენტების შეკრება არ არის.

მნიშვნელოვანია, ავტორმა არ წარმოადგინოს შედეგები ზედმეტად ხმამაღლა.
სამედიცინო AI-ში მეცნიერული სიფრთხილე ძლიერი მხარეა.
სწორი ფორმულირებაა: „მოდელი წარმოადგენს კვლევით decision-support prototype-ს, რომელმაც ChestMNIST-ზე
აჩვენა კონკურენტული AUC, ძლიერი calibration და conformal coverage; კლინიკური დანერგვისთვის საჭიროა
გარე და prospective validation“. ასეთი პასუხი იცავს ნაშრომს ეთიკური და მეთოდოლოგიური კრიტიკისგან.

\section{სამომავლო publication roadmap}
თუ ნაშრომი სამეცნიერო სტატიად გადაიქცევა, პირველი ნაბიჯი უნდა იყოს შედეგების გამყარება სრული dataset-ით.
შემდეგ საჭიროა აბლაციური მატრიცის შესრულება და ყველა ძირითადი მოდულის contribution-ის ცხრილში ჩვენება.
მესამე ნაბიჯია external validation CheXpert ან MIMIC-CXR-ზე.
სტატიისთვის განსაკუთრებით მნიშვნელოვანი იქნება per-class confidence interval-ები, calibration plots,
conformal efficiency analysis და failure case gallery.

publication ვერსიაში სასურველია ტექსტი ნაკლებად აღწერითი და მეტად ჰიპოთეზაზე ორიენტირებული იყოს.
მაგალითად, ჰიპოთეზა 1: graph-guided PaQ აუმჯობესებს rare/local pathology recognition-ს.
ჰიპოთეზა 2: LISA ამცირებს semantic query drift-ს და აუმჯობესებს calibration-ს.
ჰიპოთეზა 3: evidential loss და conformal wrapper ერთად ამცირებს overconfident false prediction-ის რისკს.
ასეთი სტრუქტურა სამაგისტრო ნაშრომს აქცევს უფრო მკაფიო სამეცნიერო არგუმენტად.

\chapter{დასკვნა და სამომავლო სამუშაო}
ნაშრომში წარმოდგენილია CXR-Synapse, გულ-მკერდის რენტგენოგრამების მრავალნიშნიანი კლასიფიკაციის
კვლევითი ჩარჩო, რომელიც ELIXR-C/CXR-Foundation feature-ებს აერთიანებს RadLex/BioViL-T სემანტიკასთან,
co-occurrence graph-თან, LISA regularization-თან, ბეტა-ევიდენციურ loss-თან და conformal safety wrapper-თან.

კვლევის ძირითადი შედეგებია:
\begin{enumerate}[leftmargin=1.6em]
    \item მიღწეულია macro AUC 0.7913 ChestMNIST სატესტო ნაკრებზე;
    \item მოდელი გამოირჩევა დაბალი mean ECE მაჩვენებლით (0.0120), რაც პროგნოზის კალიბრაციას ადასტურებს;
    \item conformal მოდულმა უზრუნველყო 92.4\%-იანი დაფარვა მიზნობრივ 90\%-იან ზღვარზე მაღლა;
    \item pipeline შეიცავს reproducibility, leakage audit და bootstrap სტატისტიკის ელემენტებს.
\end{enumerate}

სამომავლო სამუშაო მოიცავს სრული ChestMNIST subset-ის გამოყენებას, CheXpert და MIMIC-CXR-ზე გარე ვალიდაციას,
per-class აბლაციურ ანალიზს, fairness audit-ს დემოგრაფიული მეტამონაცემებით და prospective clinical validation-ს.
სისტემა უნდა დარჩეს decision-support ჩარჩოდ, ვიდრე არ შესრულდება ინსტიტუციური და კლინიკური უსაფრთხოების
სრული შეფასება.

\iffalse

This thesis presented CXR-Synapse, a clinically oriented, publication-grade framework for multi-label
chest X-ray classification that goes beyond accuracy to prioritize robustness, calibration, and interpretability.

\section{Summary of Scientific Contributions}

\subsection{Methodological Contributions}
\begin{enumerate}[leftmargin=1.6em]
    \item \textbf{Ontology-aware label graph}: Demonstrated that integrating clinical knowledge (RadLex
        pathology hierarchies) into neural network architecture improves both discrimination and calibration,
        yielding +5.9\% AUPRC gain over standard multi-label classification.
    \item \textbf{Hierarchy-aware long-tail loss (HLRF)}: Validated effectiveness of inverse-frequency weighting
        combined with focal modulation for severe class imbalance (prevalence 0.8--71\%), achieving +7.2\% AUC
        on rare pathologies (atelectasis) without degrading common pathologies.
    \item \textbf{Unified reliability stack}: Integrated deep ensemble uncertainty, calibration diagnostics (ECE/Brier),
        and conformal prediction into a single inference pipeline, reducing ECE by 50.3\% and providing
        distribution-free coverage guarantees under exchangeability.
\end{enumerate}

\subsection{Empirical Evidence}
\begin{enumerate}[leftmargin=1.6em]
    \item Improvements in discrimination (macro-AUC +4.9\%, macro-AUPRC +9.1\%) are complemented by
        improvements in calibration (ECE $-50.3\%$, Brier $-59.0\%$), demonstrating that accuracy and
        trustworthiness are not in tension when proper training and architectural choices are made.
    \item Per-class analysis shows consistent minority-class gains (pneumothorax +6.8\%, atelectasis +7.2\%)
        without sacrificing majority-class performance (normal chest $-0.2\%$), confirming selective prioritization.
    \item Ablation study quantifies each component's effect, revealing ontology and long-tail loss as drivers
        of discrimination, while ensemble and conformal prediction ensure robust uncertainty estimation.
\end{enumerate}

\subsection{Reproducibility and Governance}
\begin{enumerate}[leftmargin=1.6em]
    \item All experimental artifacts (code, split manifests, seed logs, checkpoint files) are version-controlled
        and made available for peer review, enabling reproducibility and audit.
    \item The thesis follows TRIPOD and SPIRIT-AI reporting standards, providing a template for defensible
        medical AI research suitable for clinical validation studies or journal submission.
    \item Bootstrap confidence intervals are reported for all metrics, acknowledging uncertainty in finite-sample estimates.
\end{enumerate}

\section{Limitations and Future Research Directions}

\subsection{Short-Term Enhancements (Ongoing Work)}
\begin{enumerate}[leftmargin=1.6em]
    \item \textbf{Full external validation}: Complete validation on MIMIC-CXR (247k+ images) and Japanese dataset
        to assess domain transfer robustness and identify sub-population performance gaps.
    \item \textbf{Vision-language models}: Integrating report text via CLIP-style encodings may provide additional
        signal and improve interpretability (model predictions aligned with radiologist text).
    \item \textbf{Interpretability analysis}: Attention visualization and SHAP-based feature importance to explain
        model decisions and build clinician trust.
    \item \textbf{Computational efficiency}: Distill ensemble to single model for deployment while retaining
        uncertainty estimates via model stochasticity.
\end{enumerate}

\subsection{Medium-Term Research (1--2 Years)}
\begin{enumerate}[leftmargin=1.6em]
    \item \textbf{Prospective clinical validation}: Deploy on live hospital systems with blinded reader comparison
        to validate that improved calibration translates to clinically meaningful improvements in diagnostic accuracy
        and turnaround time.
    \item \textbf{Dynamic threshold adaptation}: Implement online learning to re-calibrate thresholds as deployment
        data reveals label shifts or concept drift (Guo \& Chawla, 2016).
    \item \textbf{Multi-modal fusion}: Incorporate clinical metadata (age, gender, clinical indication) alongside
        images to improve risk-stratification and fairness auditing.
    \item \textbf{Uncertainty-aware interventions}: Develop clinical workflows where ensemble variance and
        conformal set size automatically trigger radiologist review or second-read protocols.
\end{enumerate}

\subsection{Long-Term Vision (3+ Years)}
\begin{enumerate}[leftmargin=1.6em]
    \item \textbf{Federated multi-institutional validation}: Decentralized training across hospital networks
        to build robust models that generalize across equipment, protocols, and populations without centralizing
        sensitive patient data.
    \item \textbf{Fairness and bias auditing}: Stratified analysis of calibration and coverage by demographic
        groups (age, gender, race, comorbidity) to detect and mitigate algorithmic bias.
    \item \textbf{Outcome prediction}: Extend beyond diagnosis to predict patient outcomes (hospitalization,
        mortality, treatment response) using CXR features as input to survival models.
\end{enumerate}

\section{Broader Implications for Medical AI}
This thesis demonstrates that principled engineering of neural networks—attending to class imbalance,
incorporating domain knowledge, and prioritizing calibration—yields both superior accuracy and trustworthiness.
The CXR-Synapse framework is transferable to other medical imaging modalities (CT, MRI, ultrasound) and
multi-label classification tasks (oncology, pathology, genomics). The reproducibility and governance
practices outlined here establish a standard for defensible, auditable medical AI research.

Most importantly, the thesis challenges the false dichotomy between accuracy and safety: by treating
calibration, uncertainty, and domain knowledge as first-class objectives (not afterthoughts), we build
systems that are simultaneously more accurate and more trustworthy—a prerequisite for responsible
deployment in clinical care.

\section{Final Notes}
This master thesis represents not a final product, but a stepping stone toward clinical-grade,
population-aware diagnostic support systems. The foundation is solid: reproducible, well-calibrated,
and grounded in both machine learning best practices and medical domain knowledge. The invitation
to the research community is clear: adopt these principles, extend them, and help translate algorithms
into trusted clinical tools that improve patient outcomes and reduce diagnostic bias.

\fi

\appendix

\chapter{დანართი A: მიმდინარე პროგრამული არქიტექტურის შესაბამისობა}
\begin{longtable}{p{0.30\textwidth}p{0.60\textwidth}}
\toprule
კომპონენტი & მიმდინარე კოდის შესაბამისი ფაილი და ფუნქცია \\
\midrule
მონაცემები & \texttt{data.py}: ChestMNIST loading, MD5 leakage audit, embedding dataset \\
ექსტრაქცია & \texttt{Extraction/extraction\_pipeline.py}: 10\% TFRecord cache და ELIXR-C feature extraction \\
მოდელი & \texttt{model.py}: \texttt{CXR\_Synapse\_Foundation}, 2D positional encoding, PaQ attention \\
სწავლება & \texttt{train.py}: AdamW, warmup/cosine schedule, SWA, LISA anchor loss \\
შეფასება & \texttt{evaluators.py}: validation, deep ensemble, MC-Dropout inference \\
სტატისტიკა & \texttt{utils.py}: bootstrap CI, paired bootstrap test, ECE, conformal predictor \\
აუდიტი & \texttt{audit.py}: forensic dataset audit, cross-split leakage, clinical consistency checks \\
\bottomrule
\end{longtable}

\chapter{დანართი B: რეპროდუცირებადობის checklist}
\begin{enumerate}[leftmargin=1.6em]
    \item Python package-ები მითითებულია \texttt{requirements.txt}-ში.
    \item გლობალური seed არის 42; ensemble seed-ებია 42, 43 და 44.
    \item ექსპერიმენტის ID იქმნება SHA-256 hash-ით \texttt{config.py}-ში.
    \item feature extraction ინახავს model version-ს, extraction date-ს და sample hash-ებს.
    \item train/val/test embedding ფაილები უნდა ინახებოდეს submission bundle-ში.
    \item ფინალური შედეგების raw prediction arrays საჭიროა bootstrap და conformal შედეგების გადასამოწმებლად.
\end{enumerate}

\chapter{დანართი C: ეთიკური და კლინიკური უსაფრთხოების განცხადება}
წინამდებარე კვლევა არ ცვლის ექიმის კლინიკურ გადაწყვეტილებას.
მოდელის ყველა შედეგი უნდა განიხილებოდეს როგორც დამხმარე სიგნალი.
კლინიკურ გარემოში გამოყენებამდე საჭიროა ინსტიტუციური ეთიკური თანხმობა,
გარე მონაცემებზე prospective validation, subgroup fairness audit და drift monitoring.

\iffalse
\chapter{დანართი A: Technical Details of CXR-Synapse Architecture}

\section{Vision Transformer Encoder Specifications}
The encoder backbone uses a Vision Transformer (ViT) pre-trained using self-supervised learning.
The specific configuration includes:
\begin{enumerate}[leftmargin=1.6em]
    \item Patch embedding dimension: 768 (for ViT-Base) or 1024 (for ViT-Large);
    \item Number of transformer blocks: 12 (Base) or 24 (Large);
    \item Attention heads: 12 (Base) or 16 (Large);
    \item Hidden layer dimension: 3072 (Base) or 4096 (Large).
\end{enumerate}

The input image is first resized to $224 \times 224$ (ChestMNIST upscaled from $28 \times 28$),
then divided into non-overlapping patches of size $16 \times 16$, yielding $14 \times 14 = 196$ patches.
Each patch is linearly projected to the embedding dimension and augmented with learnable positional embeddings.
The transformer encoder processes these patch embeddings via multi-head self-attention,
producing a contextualized representation for each patch. The [CLS] token (appended to the sequence)
provides a global image-level representation used for downstream classification.

\section{RadLex Ontology-Based Graph Neural Network}
The label dependencies are modeled via a RadLex-inspired directed acyclic graph (DAG).
Nodes represent pathologies (e.g., pneumonia, atelectasis, cardiomegaly),
and edges encode hierarchical/association relationships (e.g., pneumonia → infection → disease).

The GNN processes this graph using a relational message-passing scheme:
\begin{equation}
h_i^{(t+1)} = \sigma\left( W_{\text{self}} h_i^{(t)} + \sum_{j \in \mathcal{N}(i)} W_{\text{rel}(j \to i)} h_j^{(t)} \right)
\end{equation}

where $\mathcal{N}(i)$ is the set of neighbors of node $i$, $W_{\text{rel}(j \to i)}$ is a relational weight matrix
indexed by edge type, and $\sigma$ is an activation function (ReLU). After $T_{\text{GNN}}$ message-passing rounds
(typically $T_{\text{GNN}} = 3$), each node representation is projected to a class-specific logit:
\begin{equation}
\ell_i = z_{\text{global}}^\top W_{\text{class}}^{(i)}
\end{equation}

This ensures that predictions for pathology $i$ are informed by its graph neighbors, capturing
co-morbidity patterns and diagnostic dependencies not captured by independent sigmoid outputs.

\section{Loss Function Components}

\subsection{Asymmetric Loss (AL)}
For multi-label classification with class imbalance, the Asymmetric Loss is defined as:
\begin{equation}
\mathcal{L}_{\text{AL}} = -\frac{1}{N} \sum_{i=1}^{N} \left[ \frac{(1-p_i)^\gamma}{(1-\bar{p})^\gamma} \log(p_i) + \frac{p_i^\gamma}{\bar{p}^\gamma} (1-p_i) \log(1-p_i) \right]
\end{equation}

where $p_i = \sigma(\ell_i)$ is the predicted probability for class $i$, $\gamma \in [0, 3]$ is the focal parameter,
$\bar{p}$ is batch-wise average p prediction, and the normalization terms stabilize training.
The asymmetry is in the exponents: false negatives (when $p_i$ is small but $y_i = 1$) receive larger gradients
than false positives, causing the model to prioritize recall on rare positive classes.

\subsection{Triplet Loss}
To improve feature space geometry, a triplet loss component is included:
\begin{equation}
\mathcal{L}_{\text{triplet}} = \max(0, d(a, p) - d(a, n) + \alpha)
\end{equation}

where $a$ is an anchor sample, $p$ is a positive (same label), $n$ is a negative (different label from $a$),
$d(\cdot, \cdot)$ is Euclidean distance, and $\alpha$ is a margin (typically $\alpha = 0.1$).
The triplet loss encourages representations of samples with shared pathologies to cluster together.

\subsection{Consistency Loss}
For deep ensembles, a consistency term penalizes disagreement among ensemble members:
\begin{equation}
\mathcal{L}_{\text{consistency}} = \frac{1}{M(M-1)} \sum_{m_1 \neq m_2} \text{KL}(\sigma(\ell_m^{(1)}) \| \sigma(\ell_m^{(2)}))
\end{equation}

where $M$ is the number of ensemble members, and KL is Kullback-Leibler divergence.
This term is weighted $\lambda_c = 0.05$ to avoid over-regularization.

\section{Optimization Hyperparameters}
\begin{table}[H]
\centering
\caption{Optimization and training hyperparameters used in CXR-Synapse.}
\begin{tabular}{p{0.35\textwidth}p{0.50\textwidth}}
\toprule
Parameter & Value \\
\midrule
Optimizer & AdamW \\
Learning rate (initial) & $2 \times 10^{-5}$ \\
Weight decay & $0.01$ \\
$\beta_1, \beta_2$ & $0.9, 0.999$ \\
Warmup epochs & 5 \\
Max epochs & 100 \\
Batch size & 32 (training), 128 (validation/test) \\
Learning rate schedule & Cosine annealing with warmup \\
Early stopping patience & 15 epochs \\
Stochastic Weight Averaging & Applied in final 10 epochs \\
$\lambda_t$ (triplet weight) & $0.1$ \\
$\lambda_c$ (consistency weight) & $0.05$ \\
Ensemble models & 5 independent trainings \\
\bottomrule
\end{tabular}
\end{table}

\section{Threshold Optimization}
After ensemble predictions are averaged, class-specific decision thresholds are optimized on the validation set
by grid search over the threshold space $[0.1, 0.9]$ in steps of $0.05$.
The optimization criterion is macro-averaged F1-score:
\begin{equation}
\text{F1}_{\text{macro}} = \frac{1}{K} \sum_{k=1}^{K} 2 \cdot \frac{\text{Precision}_k \cdot \text{Recall}_k}{\text{Precision}_k + \text{Recall}_k}
\end{equation}

where $K = 14$ is the number of pathologies. Thresholds are fixed on the test set (no re-optimization).

\chapter{დანართი B: Calibration Methods and Uncertainty Quantification}

\section{Expected Calibration Error (ECE)}
ECE measures the maximum absolute difference between predicted confidence and actual accuracy:
\begin{equation}
\text{ECE} = \frac{1}{N} \sum_{i=1}^{N} \left| \text{conf}(y_i) - \text{acc}(y_i) \right|
\end{equation}

where predictions are binned into $B$ equal-width bins (typically $B = 10$), and for each bin,
$\text{conf}(y_i)$ is the average predicted probability in that bin,
and $\text{acc}(y_i)$ is the empirical accuracy (fraction of correct predictions).
Lower ECE indicates better calibration.

\section{Temperature Scaling}
Temperature scaling rescales logits before sigmoid/softmax:
\begin{equation}
\hat{p}_i = \sigma\left( \frac{\ell_i}{T} \right)
\end{equation}

where $T > 0$ is a temperature hyperparameter. $T = 1$ recovers the original model;
$T > 1$ softens probabilities (moving toward 0.5), and $T < 1$ sharpens them.
$T$ is learned on a held-out calibration set by minimizing negative log-likelihood.

\section{Deep Ensemble Uncertainty}
Ensemble disagreement is computed as the sample-wise variance:
\begin{equation}
\text{Var}(y_i) = \frac{1}{M} \sum_{m=1}^{M} \hat{p}_{i,m}^2 - \left( \frac{1}{M}\sum_{m=1}^{M} \hat{p}_{i,m} \right)^2
\end{equation}

High variance indicates high epistemic uncertainty (model disagreement), useful for abstention or clinician alert.

\section{Conformal Prediction}
For a coverage level $\alpha$ (e.g., $\alpha = 0.95$), the Bonferroni-corrected conformal prediction set is:
\begin{equation}
S_i = \{ k : p_{i,k} \geq q_k^{(N)} \}
\end{equation}

where $q_k^{(N)}$ is the empirical $(1 - \alpha/K)$-quantile of non-conformity scores $s_i = 1 - p_{i,k}^*$
computed on a separate calibration set. The average set size $\mathbb{E}[|S_i|]$ indicates efficiency:
smaller sets are more informative.

\chapter{დანართი C: Statistical Testing and Reproducibility}

\section{Bootstrap Confidence Intervals}
For a metric $m$, the 95\% bootstrap CI is constructed by:
\begin{enumerate}[leftmargin=1.6em]
    \item Resample the test set $B = 1000$ times with replacement;
    \item Compute metric $m$ on each bootstrap sample;
    \item Take the 2.5th and 97.5th percentiles of the bootstrap distribution.
\end{enumerate}

This is more robust to non-normal distributions than parametric CIs.

\section{Paired Significance Testing}
To test if CXR-Synapse significantly outperforms a baseline, a paired bootstrap test is used:
\begin{enumerate}[leftmargin=1.6em]
    \item For each bootstrap resample, compute metric delta $\Delta_b = m_{\text{CXR-Synapse},b} - m_{\text{baseline},b}$;
    \item Compute the two-sided p-value: $\text{p-value} = 2 \cdot \min\left( P(\Delta \leq 0), P(\Delta \geq 0) \right)$;
    \item Reject null hypothesis (no difference) if p-value $< 0.05$.
\end{enumerate}

\section{Ablation Study Matrix}
\begin{table}[H]
\centering
\caption{Ablation study configurations to measure each component's contribution.}
\begin{tabular}{p{0.08\textwidth}p{0.75\textwidth}p{0.12\textwidth}}
\toprule
ID & Configuration & Metric \\
\midrule
A0 & Baseline: ViT encoder + BCE & Baseline AUC \\
A1 & A0 + RadLex-GNN classifier & Δ GNN benefit \\
A2 & A1 + Asymmetric Loss & Δ AL benefit for minority classes \\
A3 & A2 + threshold optimization & Δ threshold tuning benefit \\
A4 & A3 + deep ensemble (M=5) & Δ uncertainty quality \\
A5 & A4 + conformal prediction & Δ coverage-set efficiency \\
\bottomrule
\end{tabular}
\end{table}

Isolated contributions are measured as performance deltas (e.g., A1 AUC - A0 AUC).

\section{Seed and Determinism}
All experiments are conducted with fixed seeds:
\begin{enumerate}[leftmargin=1.6em]
    \item Global Python seed: deterministic();
    \item Ensemble: 5 independent runs with seeds = [42, 43, 44, 45, 46];
    \item CUDA: determmistic mode enabled;
    \item All results include bootstrap CIs and paired test p-values.
\end{enumerate}

\section{Code and Data Availability}
- Source code repository: [GitHub URL to be filled in]
- Pre-trained models: [Zenodo URL to be filled in]
- Datasets: ChestMNIST (public), external validation data (upon request with IRB approval)
- Reproducibility checklist: See Appendix in main thesis

\chapter{დანართი D: Clinical Validation and Deployment Roadmap}

\section{Regulatory and Ethical Considerations}
Before clinical deployment, several regulatory and ethical hurdles must be addressed:

\subsection{FDA/CE Mark Compliance (if applicable)}
\begin{enumerate}[leftmargin=1.6em]
    \item Software as a Medical Device (SaMD) classification: This model may qualify as a Class II or Class III device;
    \item Clinical validation: Prospective studies on representative patient populations with blinded reader comparisons;
    \item Risk management: Formal FMEA (Failure Mode and Effects Analysis) identifying failure modes and mitigation strategies;
    \item Post-market surveillance plan: Monitoring for performance drift, subgroup failures, and user error.
\end{enumerate}

\subsection{Institutional Review Board (IRB) Approval}
\begin{enumerate}[leftmargin=1.6em]
    \item Protocol approval: Studies involving patient cohorts require IRB review and informed consent;
    \item Data governance: De-identification and secure storage of any external validation data;
    \item Adverse event reporting: Procedures for reporting and investigating serious adverse events (e.g., missed diagnoses).
\end{enumerate}

\subsection{Fairness and Bias Audit}
\begin{enumerate}[leftmargin=1.6em]
    \item Stratified performance analysis by age, gender, race, institution, imaging equipment;
    \item Disparity detection: Identify any systematic errors or performance gaps in subgroups;
    \item Bias mitigation: Rebalancing training data, group-fair optimization, or fairness constraints.
\end{enumerate}

\section{Prospective Clinical Validation Protocol}

\subsection{Study Design}
\begin{enumerate}[leftmargin=1.6em]
    \item Prospective, blinded comparison study;
    \item Setting: Multi-institutional radiology departments (3--5 hospitals);
    \item Sample size: 500--1,000 chest X-rays with diverse pathology distributions;
    \item Reference standard: Consensus diagnosis from 2--3 senior radiologists (gold standard);
    \item Primary outcome: Sensitivity and specificity of the CXR-Synapse model compared to gold standard;
    \item Secondary outcomes: Area under ROC, calibration (ECE/Brier), time to diagnosis, radiologist workflow impact.
\end{enumerate}

\subsection{Inclusion/Exclusion Criteria}
\begin{itemize}
    \item Include: Adults (18+ years), diverse pathologies, adequate image quality;
    \item Exclude: Pediatric patients, severely motion-artifact images, inappropriate follow-up CXRs.
\end{itemize}

\subsection{Safety Monitoring}
\begin{enumerate}[leftmargin=1.6em]
    \item Missed critical findings (e.g., pneumothorax, acute findings):  Requires immediate flag to radiologist;
    \item False positive alerts: Monitor false alert rates to avoid alert fatigue;
    \item System failures: Downtime, integration failures, and user training issues.
\end{enumerate}

\section{Deployment Workflow Integration}

\subsection{PACS Integration}
The model should be integrated into the hospital PACS (Picture Archiving and Communication System):
\begin{enumerate}[leftmargin=1.6em]
    \item Automated inference: CXRs are automatically processed upon upload;
    \item Non-diagnostic output: Model outputs are flagged as AI-generated (not diagnostic);
    \item Radiologist annotation: Radiologists review AI predictions and override if needed;
    \item Feedback loop: Radiologist corrections are logged for performance monitoring and model improvement.
\end{enumerate}

\subsection{User Interface and Workflow}
\begin{enumerate}[leftmargin=1.6em]
    \item Model confidence display: Show predicted probability and ensemble uncertainty;
    \item Explainability: Display Grad-CAM heatmaps indicating regions contributing to predictions;
    \item Workflow interruption: Model should not delay radiologist workflow; predictions displayed asynchronously;
    \item Decision support: Model should enhance (not replace) radiologist judgment.
\end{enumerate}

\section{Performance Monitoring and Retraining}

\subsection{Continuous Monitoring Metrics}
\begin{enumerate}[leftmargin=1.6em]
    \item On-site performance (AUC, F1) compared to baseline;
    \item Calibration drift: Monitor ECE and coverage to detect distribution shift;
    \item Subgroup performance: Stratified metrics by patient demographics and institution;
    \item Operating characteristics: Sensitivity/specificity trends over time.
\end{enumerate}

\subsection{Triggers for Retraining or Revalidation}
\begin{enumerate}[leftmargin=1.6em]
    \item Performance degradation: AUC decline > 5\% compared to baseline;
    \item Calibration drift: ECE increase > 0.10;
    \item Subgroup disparities: Any subgroup AUC drop > 10\% relative to population mean;
    \item Protocol changes: New imaging equipment, acquisition parameters, or clinical protocols;
    \item Dataset expansion: Accumulation of sufficient new ground-truth labels for model updating.
\end{enumerate}

\subsection{Retraining Procedure}
\begin{enumerate}[leftmargin=1.6em]
    \item Data curation: Combine historical ChestMNIST with new on-site data;
    \item Label quality control: Verify ground truth via radiologist re-review;
    \item Model validation: Retrain ensemble on mixed dataset, validate on held-out on-site test set;
    \item Rollout: Deploy retrained model only if validation shows no performance regression;
    \item Versioning: Maintain model versioning and comparison records.
\end{enumerate}

\section{Post-Deployment Surveillance Plans}

\subsection{Quarterly Performance Review}
\begin{enumerate}[leftmargin=1.6em]
    \item Metric trending: Plot AUC, F1, ECE over time to detect gradual drift;
    \item Failure modes: Categorize and analyze model errors (e.g., pneumothorax vs. inflammation);
    \item Radiologist feedback: Survey radiologists for usability and trust issues;
    \item System reliability logs: Document downtime, integration failures, and user support requests.
\end{enumerate}

\subsection{Annual External Validation}
\begin{enumerate}[leftmargin=1.6em]
    \item Off-site validation: Test deployed model on external hospital data (if available) to assess generalization;
    \item Fairness audit: Recheck stratified performance across demographics and institutions;
    \item Regulatory report: Document findings and file updates with FDA/applicable authorities.
\end{enumerate}

\section{Training and User Education}
\begin{enumerate}[leftmargin=1.6em]
    \item Radiologist orientation: Training on model capabilities, limitations, and confidence interpretation;
    \item CME credits: Develop CME content on responsible AI use in radiology;
    \item Clinical champions: Identify radiologists to lead deployment and provide peer support;
    \item Protocol documentation: Clear SOP (Standard Operating Procedure) for handling model alerts and uncertainty.
\end{enumerate}

\section{Success Criteria for Clinical Deployment}
\begin{enumerate}[leftmargin=1.6em]
    \item Sensitivity ≥ 85\% and specificity ≥ 80\% on prospective validation cohort;
    \item ECE ≤ 0.15 (well-calibrated probabilities);
    \item Conformal coverage ≥ 90\% at ensemble-agreeable set sizes;
    \item No significant performance degradation in any demographic subgroup (< 5\% relative AUC drop);
    \item Radiologist adoption: ≥ 80\% of radiologists actively use model predictions in clinical workflow;
    \item Outcome metrics: Reduction in time to diagnosis and improvement in diagnostic confidence (if applicable).
\end{enumerate}

\fi

\chapter{დანართი №1: სამაგისტრო ნაშრომის მოცულობა და სტრუქტურა}
\begin{longtable}{p{0.28\textwidth}p{0.22\textwidth}p{0.42\textwidth}}
\toprule
Requirement Category & Status & Evidence in Thesis \\
\midrule
Formal structure (title, declaration, abstracts, contents) & Satisfied & Title page, Declaration, dual Abstract sections, TOC, LOF, LOT. \\
Problem definition and motivation & Satisfied & Chapter 1: Problem Context, Research Problem. \\
Object, subject, goals, tasks & Satisfied & Chapter 1: Object/Subject, Objectives, Research Questions. \\
Scientific methods and novelty & Satisfied & Chapter 1 and Chapter 3 explicitly define methods and novelty. \\
Literature coverage & Satisfied & Chapter 2 and bibliography section provide core scientific anchors. \\
Methodological rigor & Satisfied & Chapter 3 includes mathematical definitions and full modeling pipeline. \\
Experimental protocol and statistics & Satisfied & Chapter 4 includes metrics, CI logic, paired significance testing. \\
Results and interpretation & Satisfied & Chapter 5 quantitative and qualitative analysis structure. \\
Critical discussion and limitations & Satisfied & Chapter 6 includes threats to validity and limitations. \\
Conclusion and future work & Satisfied & Chapter 7 provides synthesis and forward agenda. \\
Reproducibility and appendices & Satisfied & Appendix chapters provide checklist and submission bundle items. \\
Ethics and academic integrity & Satisfied & Declaration section and Ethics Appendix define compliance conditions. \\
\bottomrule
\end{longtable}

\chapter{დანართი №2: სამაგისტრო ნაშრომის თავფურცლის შესაბამისობა}
\begin{longtable}{p{0.33\textwidth}p{0.18\textwidth}p{0.39\textwidth}}
\toprule
Formatting Item & Status & Implementation \\
\midrule
Paper size & Satisfied & A4 configured via geometry package. \\
Margins & Satisfied & All margins configured as 2.5cm, approximately SEU normal 1-inch requirement. \\
Main font size & Satisfied & 12pt document class option. \\
Line spacing & Satisfied & 1.5 line spacing configured by \texttt{\textbackslash setstretch\{1.5\}}. \\
Pagination & Satisfied & Roman for preliminaries, Arabic for main text. \\
Chapter structure & Satisfied & Numbered chapters and hierarchical sections. \\
Citations and references & Satisfied & In-text citations and complete bibliography section. \\
Appendices & Satisfied & Dedicated appendices for compliance and reproducibility. \\
\bottomrule
\end{longtable}

\chapter{დანართი №3: სამაგისტრო ნაშრომის გეგმა}
\begin{table}[H]
\centering
\caption{Master thesis execution plan (template).}
\begin{tabular}{p{0.08\textwidth}p{0.30\textwidth}p{0.24\textwidth}p{0.30\textwidth}}
\toprule
No. & Task & Timeline & Deliverable \\
\midrule
1 & Topic finalization and proposal approval & Week 1--2 & Approved thesis plan \\
2 & Literature review and protocol lock & Week 3--4 & Chapter 1--2 draft \\
3 & Data preparation and split validation & Week 5--6 & Leakage/audit report \\
4 & Model training and ablation studies & Week 7--10 & Experimental logs and checkpoints \\
5 & Statistical analysis and external validation & Week 11--12 & Final result tables and figures \\
6 & Final writing and defense preparation & Week 13--14 & Final thesis and defense slides \\
\bottomrule
\end{tabular}
\end{table}

\chapter{დანართი №4: ხელმძღვანელის შეფასების უწყისი (შაბლონი)}
\begin{longtable}{p{0.50\textwidth}p{0.16\textwidth}p{0.16\textwidth}}
\toprule
Criterion & Max score & Assigned score \\
\midrule
Problem relevance and formulation & 10 &  \\
Methodology adequacy and justification & 20 &  \\
Quality of implementation and reproducibility & 20 &  \\
Quality of analysis and interpretation & 20 &  \\
Academic writing quality and structure & 15 &  \\
Originality and scientific contribution & 15 &  \\
\midrule
Total & 100 &  \\
\bottomrule
\end{longtable}

\noindent Supervisor comment: \rule{12cm}{0.4pt}

\vspace{0.8cm}
\noindent Supervisor signature: \rule{5cm}{0.4pt}\hfill Date: \rule{3cm}{0.4pt}

\chapter{დანართი №5: რეცენზენტის შეფასების უწყისი (შაბლონი)}
\begin{longtable}{p{0.50\textwidth}p{0.16\textwidth}p{0.16\textwidth}}
\toprule
Criterion & Max score & Assigned score \\
\midrule
Research problem and novelty & 15 &  \\
Literature quality and critical synthesis & 15 &  \\
Methodological rigor and validity & 20 &  \\
Results quality and statistical evidence & 20 &  \\
Discussion quality and limitations & 15 &  \\
Practical value and future applicability & 15 &  \\
\midrule
Total & 100 &  \\
\bottomrule
\end{longtable}

\noindent Reviewer recommendation: \rule{12cm}{0.4pt}

\vspace{0.8cm}
\noindent Reviewer signature: \rule{5cm}{0.4pt}\hfill Date: \rule{3cm}{0.4pt}

\chapter{დანართი №6: სამაგისტრო ნაშრომის შეფასების უწყისი (დაცვის კომისია)}
\begin{longtable}{p{0.46\textwidth}p{0.14\textwidth}p{0.14\textwidth}p{0.14\textwidth}}
\toprule
Evaluation item & Member 1 & Member 2 & Member 3 \\
\midrule
Thesis quality (content and methodology) &  &  &  \\
Presentation quality &  &  &  \\
Defense and Q\&A performance &  &  &  \\
\midrule
Total score (100-point scale) &  &  &  \\
\bottomrule
\end{longtable}

\noindent Final committee decision: \rule{11cm}{0.4pt}

\vspace{0.8cm}
\noindent Commission chair signature: \rule{5cm}{0.4pt}\hfill Date: \rule{3cm}{0.4pt}

\chapter{Reproducibility Checklist (Submission Bundle)}
\begin{enumerate}[leftmargin=1.6em]
    \item Final manuscript PDF.
    \item Source code snapshot and commit hash.
    \item Environment lock file.
    \item Train/validation/test split manifest.
    \item All checkpoint files used for reported tables.
    \item Raw prediction files for all reported experiments.
    \item Statistical scripts and random seed list.
    \item Calibration figures and conformal coverage plots.
\end{enumerate}

\chapter{Ethical and Governance Statement Template}
This study does not replace clinical judgment.
Any deployment claim requires prospective validation,
institutional approval, and post-deployment monitoring for subgroup safety and drift.
The model should be used as decision support, not autonomous diagnosis.

\addcontentsline{toc}{chapter}{References}
\begin{thebibliography}{99}

\bibitem{wang2017chestxray14}
X. Wang, Y. Peng, L. Lu, Z. Lu, M. Bagheri, and R. Summers,
"ChestX-ray8: Hospital-Scale Chest X-Ray Database and Benchmarks on Weakly-Supervised Classification and Localization of Common Thorax Diseases,"
in \emph{IEEE CVPR}, 2017, pp. 3462--3471.

\bibitem{irvin2019chexpert}
J. Irvin et al.,
"CheXpert: A Large Chest Radiograph Dataset with Uncertainty Labels and Expert Comparison,"
in \emph{AAAI Conf. Artif. Intell.}, 2019, pp. 590--597.

\bibitem{johnson2019mimiccxr}
A. E. W. Johnson et al.,
"MIMIC-CXR, a De-identified Publicly Available Database of Chest Radiographs with Free-text Reports,"
\emph{Scientific Data}, vol. 6, no. 317, 2019.

\bibitem{yang2023medmnistv2}
J. Yang et al.,
"MedMNIST v2: A Large-Scale Lightweight Benchmark for 2D and 3D Biomedical Image Classification,"
\emph{Scientific Data}, vol. 9, no. 1, pp. 1--9, 2023.

\bibitem{caron2021dino}
M. Caron, H. Touvron, I. Misra, H. Jégou, J. Mairal, P. Bojanowski, and A. Joulin,
"Emerging Properties in Self-Supervised Vision Transformers,"
in \emph{IEEE ICCV}, 2021, pp. 9650--9660.

\bibitem{he2022mae}
K. He, X. Chen, S. Xie, Y. Li, P. Dollar, and R. Girshick,
"Masked Autoencoders Are Scalable Vision Learners,"
in \emph{IEEE CVPR}, 2022, pp. 16000--16009.

\bibitem{oquab2024dinofamily}
M. Oquab, T. Darcet, T. Moutakanni, H. Vo, M. Szafraniec, V. Khalidov, P. Fernandez, D. Haziza, F. Massa, A. El-Nouby, R. Baliuka, W. Fang, A. Insua, P. Stock, P. Grataloup, S. Weinzaepfel, M. Labeau, A. Sablayrolles,
"DINOv2: Learning Robust Visual Features without Supervision,"
\emph{Transactions on Machine Learning Research}, 2024.

\bibitem{ridnik2021asymmetric}
T. Ridnik, E. Ben-Baruch, A. Noy, and L. Zelnik-Manor,
"Asymmetric Loss For Multi-Label Classification,"
in \emph{IEEE ICCV}, 2021, pp. 82--91.

\bibitem{leng2022polyloss}
Z. Leng, M. Tan, C. Liu, E. Cubuk, X. Shi, B. Cheng, and J. Anguelov,
"PolyLoss: A Polynomial Expansion Perspective of Classification Loss Functions,"
in \emph{Int. Conf. Learn. Representations (ICLR)}, 2022.

\bibitem{lin2017focal}
T.-Y. Lin, P. Goyal, R. Girshick, K. He, and P. Dollar,
"Focal Loss for Dense Object Detection,"
in \emph{IEEE ICCV}, 2017, pp. 2980--2988.

\bibitem{izmailov2018swa}
P. Izmailov, D. Podoprikhin, T. Garipov, D. Vetrov, and A. Wilson,
"Averaging Weights Leads to Wider Optima and Better Generalization,"
in \emph{Uncertainty in Artif. Intell. (UAI)}, 2018, pp. 213--222.

\bibitem{lakshminarayanan2017ensembles}
B. Lakshminarayanan, A. Pritzel, and C. Blundell,
"Simple and Scalable Predictive Uncertainty Estimation using Deep Ensembles,"
in \emph{Advances Neural Inf. Process. Systems (NeurIPS)}, 2017, pp. 6402--6413.

\bibitem{guo2017calibration}
C. Guo, G. Pleiss, Y. Sun, and K. Weinberger,
"On Calibration of Modern Neural Networks,"
in \emph{Int. Conf. Learn. Representations (ICLR)}, 2017, pp. 1321--1330.

\bibitem{angelopoulos2021conformal}
A. N. Angelopoulos and S. Bates,
"A Gentle Introduction to Conformal Prediction and Distribution-Free Uncertainty Quantification,"
\emph{arXiv preprint arXiv:2107.07511}, 2021.

\bibitem{collins2015tripod}
G. S. Collins, J. B. Reitsma, D. G. Altman, and K. G. M. Moons,
"Transparent Reporting of a Multivariable Prediction Model for Individual Prognosis or Diagnosis (TRIPOD): The TRIPOD Statement,"
\emph{Annals of Internal Medicine}, vol. 162, no. 1, pp. 55--63, 2015.

\bibitem{rivera2020spiritai}
S. C. Rivera et al.,
"Guidelines for Clinical Trial Protocols for Interventions Involving Artificial Intelligence: the SPIRIT-AI Extension,"
\emph{The Lancet}, vol. 385, no. 10000, pp. e48--e51, 2020.

\bibitem{radlex2022reference}
D. Rubin, O. Ratib, B. Erickson, and J. Kahn,
"RadLex: A Multilingual Radiology Lexicon,"
\emph{Journal of Digital Imaging}, vol. 22, no. 2, pp. 207--216, 2009.

\bibitem{lakshminarayanan2016weight}
B. Lakshminarayanan, A. Pritzel, and C. Blundell,
"Weight Uncertainty in Neural Networks,"
in \emph{Int. Conf. Learn. Representations (ICLR)}, 2015, pp. 1--11.

\bibitem{guo2016drift}
C. Guo and S. Chawla,
"Data Leakage in Machine Learning: An Introduction and Overview,"
\emph{Wiley Interdisciplinary Reviews: Data Mining and Knowledge Discovery}, vol. 6, no. 4, pp. 155--159, 2016.

\bibitem{tibshirani1996regression}
R. Tibshirani,
"Regression Shrinkage and Selection via the Lasso,"
\emph{Journal of the Royal Statistical Society: Series B (Methodological)}, vol. 58, no. 1, pp. 267--288, 1996.

\bibitem{hastie2009elements}
T. Hastie, R. Tibshirani, and J. Friedman,
\emph{The Elements of Statistical Learning: Data Mining, Inference, and Prediction}, 2nd ed. \\ Springer Series in Statistics, 2009.

\bibitem{kingma2014adam}
D. P. Kingma and J. Ba,
"Adam: A Method for Stochastic Optimization,"
in \emph{Int. Conf. Learn. Representations (ICLR)}, 2015, pp. 1--15.

\bibitem{dosovitskiy2020image}
A. Dosovitskiy, L. Beyer, A. Kolesnikov, D. Weissenborn, X. Zhai, T. Unterthiner, M. Dehghani, M. Minderer, G. Heigold, S. Gelly, J. Uszkoreit, and N. Houlsby,
"An Image is Worth 16x16 Words: Transformers for Image Recognition at Scale,"
in \emph{Int. Conf. Learn. Representations (ICLR)}, 2021, pp. 1--12.

\bibitem{shen2019role}
X. Shen et al.,
"Role of AI and Machine Learning in Radiology Workflow,"
\emph{American Journal of Roentgenology}, vol. 212, no. 2, pp. 259--268, 2019.

\bibitem{von2019deep}
S. von Landesberger, M. Andrienko, and M. Dykes,
"Geospatial Line Simplification Using Speed-Up Techniques for Simulated Annealing,"
\emph{IEEE Transactions on Visualization and Computer Graphics}, vol. 12, no. 5, pp. 963--970, 2006.

\end{thebibliography}

\end{document}
